\chapter{空间向量代数与 \texorpdfstring{$O(n)$}{O(n)} 动力学递推算法（RNEA·CRBA·ABA）}
\label{ch:spatial-dynamics}

% ════════════════════════════════════════════════════════════════
%  控制部 C2 首章（卷五《机器人最优控制》刚体动力学子部之首）。
%  吸收 Robotics_Tutorial 50/10(空间向量代数 2130 行) + 50/30(O(n)递推算法 2078 行)，研究生密度。
%  一条「6D 代数工具 → O(n) 递推算法」的连续叙事线，是 K(几何力学)/L(约束·解析微分) 的前置。
%  全章空间向量用 Featherstone [omega;v] 约定；§J.2/§J.8 各放一次与本书 se(3) ξ=[rho;φ] 的置换对照。
%  右扰动主线（接口对齐 \cref{ch:lie}）；勘误 E1(Hollerbach 1980,实为O(n))–E5 已应用（设计稿 §2）。
% ════════════════════════════════════════════════════════════════

\begin{note}
\textbf{本章定位与记号约定.}\;
本章是控制部刚体动力学子部的首章。导论 \cref{ch:control} 已给出计算力矩控制、阻抗与全身控制的全景，最优控制子部 \cref{ch:opt-variational-pmp,ch:dp-hjb,ch:lqr-lqg-riccati} 已把动力学方程 $\mathbf M(\mathbf q)\ddot{\mathbf q}+\mathbf C(\mathbf q,\dot{\mathbf q})\dot{\mathbf q}+\mathbf g(\mathbf q)=\boldsymbol\tau$ 当作\emph{已知对象}来设计控制律。本章要回答一个被它们回避的计算问题：给定一棵 $n$ 自由度的关节树，如何以\textbf{与自由度数成线性}的代价计算 $\boldsymbol\tau$（逆动力学）、$\mathbf M(\mathbf q)$（质量矩阵）与 $\ddot{\mathbf q}$（正动力学）？答案是 Featherstone 的三大递推算法 RNEA/CRBA/ABA\cite{featherstone2008rigid}，它们是机械臂部、腿足部、控制部 DDP/iLQR（\cref{ch:ddp-ilqr}）与采样 MPC（\cref{ch:mpc-advanced}）共用的计算地基。

\textbf{空间向量分量序：本章统一采用 Featherstone 约定 $\hat{\mathbf v}=(\boldsymbol\omega;\mathbf v)$（角量在前、线量在后）}，与 Carpentier--Mansard\cite{carpentier2018analytical}、Lynch--Park\cite{lynch2017modern} 及主流库内部一致。\textbf{注意这与本书 $\mathfrak{se}(3)$ 坐标 $\boldsymbol\xi=[\boldsymbol\rho;\boldsymbol\phi]$（线量 $\boldsymbol\rho$ 在前）的排序相反}（见 \cref{ch:lie}），二者由置换矩阵互换；凡与李群章接口处（\cref{sec:sva-lie}）逐式标注该置换，否则跨章会把 $\mathrm{ad}$ 算子的上下三角块搞反。6×6 空间惯性记花体 $\mathcal I$（避免与单位阵 $\mathbf I$、3×3 转动惯量 $\mathbf I_c$ 混）；关节空间质量矩阵记 $\mathbf M(\mathbf q)$（$n\times n$）；变换方向 ${}^i\mathbf X_j$ 表「从坐标系 $j$ 到坐标系 $i$」。$\mathbf M,\mathbf C,\mathbf g$ 的变分来源、对称/正定/$\dot{\mathbf M}-2\mathbf C$ 反对称等性质主放几何力学章（\cref{ch:geometric-mechanics}），本章直接\emph{用}结论。
\end{note}

\section{动机与旋量理论简史}
\label{sec:sva-motivation}

\paragraph{动机：三维记号的「代数海洋」。}
考虑一个看似平凡的问题：计算旋转刚体上某点 $P$ 的速度，$\mathbf v_P=\mathbf v_O+\boldsymbol\omega\times\mathbf r_{OP}$。单独看很简单。但当处理 $n$ 自由度的串联机械臂时，每个连杆的速度都要写成前一连杆速度加相对运动，每个连杆需\emph{两个}方程（线速度与角速度）并处理叉积；进一步计算加速度时，又冒出 Coriolis 项 $\boldsymbol\omega\times(\boldsymbol\omega\times\mathbf r)$ 与离心项。对 $n$ 个连杆，加速度传递公式里的交叉耦合项呈爆炸式增长——这就是 Featherstone 所谓的\emph{代数海洋}：用三维向量写多体动力学，公式长度与复杂度随自由度急剧膨胀。

\paragraph{承诺：一个 6D 对象 + 正确的变换规则。}
空间向量框架的核心主张是：把线运动与角运动\emph{合并}成一个 6 维对象，配上一套\emph{正确的}变换、叉积与惯性规则，多体动力学的递推算法将变得极其简洁。对一个 7 自由度臂做逆动力学，三维记号需约 42 个含叉积的向量方程，而空间向量框架只需 21 个 6D 方程，无需分别处理线/角量。这不是事后的记号美化——后文将看到，递推 Newton--Euler 算法在此框架下得到统一而优雅的推导，且每一步都能映射为标准线性代数运算，便于求解析导数。

\begin{insight}
空间向量\textbf{不是}「把 $\boldsymbol\omega$ 与 $\mathbf v$ 拼成长向量」的记号便利。它的威力在于发现了 6D 空间中的\emph{代数闭合性}：所有运算（坐标变换、叉积、惯性乘法）都在同一个 6D 框架内闭合，不需要拆回 3D 做中间计算。更关键的是它改变了加速度的定义——空间加速度是空间速度的简单时间导数，不含 Coriolis 项（\cref{sec:sva-accel}），这使加速度像速度一样以加法叠加，从根本上消除了三维记号的交叉耦合。
\end{insight}

\paragraph{历史：从旋量理论到 Featherstone。}
空间向量的根源是 19 世纪的旋量理论（screw theory）。刚体一般位移可分解为「沿某轴旋转 + 沿同轴平移」的螺旋运动，这一定理的首个证明归 Giulio Mozzi（1763），意大利文献称螺旋轴为 \emph{asse di Mozzi}；Chasles（1830）是后来独立的同类工作，故定理正名为 \textbf{Mozzi--Chasles 定理}\cite{ceccarelli2000screw}（教科书常以 Chasles 简称之）。Plücker 约 1846 年引入 6 维 dyname（前 3 力、后 3 偶）的对偶向量思想，其线坐标系统在 1865 年前后成熟；Ball、Study 在 1870--1900 年代将旋量理论系统化。Yang 与 Freudenstein（\textbf{1964}）将对偶数四元数/螺旋代数用于空间机构分析\cite{yang1964application}；Murray、Li 与 Sastry（1994）从李群视角统一了 twist/wrench 理论\cite{murray1994mathematical}。

Featherstone 的空间向量框架成形于其 1984 年 Edinburgh 博士论文 \emph{Robot Dynamics Algorithms}，1987 年由 Kluwer 出版为同名专著\cite{featherstone1987robot}，2008 年扩充改写为 Springer 版 \emph{Rigid Body Dynamics Algorithms}\cite{featherstone2008rigid}（书名从 \emph{Robot} Dynamics 变为 \emph{Rigid Body} Dynamics，本书引用以 2008 版为准）；其面向初学者的直觉性导引见 IEEE RAM 2010 两部分综述\cite{featherstone2010beginner}。\cref{fig:sva-timeline} 把这条脉络绘成时间线。

\begin{figure}[ht]
\centering
\begin{tikzpicture}[>=Stealth, font=\small]
  \draw[thick,->] (0,0) -- (13.2,0) node[right]{年代};
  \foreach \x/\yr in {0.6/1763, 2.4/1846, 4.2/1870s, 6.0/1964, 7.8/1984, 9.6/2008}{
    \draw[thick] (\x,-0.1) -- (\x,0.1) node[below=3pt,font=\scriptsize]{\yr};
  }
  \node[align=center, anchor=south, font=\scriptsize] at (0.6,0.15) {Mozzi\\螺旋分解};
  \node[align=center, anchor=south, font=\scriptsize] at (2.4,0.85) {Plücker\\dyname/线坐标};
  \node[align=center, anchor=south, font=\scriptsize] at (4.2,0.15) {Ball, Study\\旋量理论系统化};
  \node[align=center, anchor=south, font=\scriptsize] at (6.0,0.85) {Yang--Freudenstein\\对偶四元数};
  \node[align=center, anchor=south, font=\scriptsize] at (7.8,0.15) {Featherstone\\6D 空间向量};
  \node[align=center, anchor=south, font=\scriptsize] at (9.6,0.85) {RBDA (2008)\\算法标准};
\end{tikzpicture}
\caption{空间向量代数的历史脉络：从 Mozzi 的螺旋分解到 Featherstone 的 6D 框架。}
\label{fig:sva-timeline}
\end{figure}

\paragraph{启下。}
下面四节（\cref{sec:sva-plucker,sec:sva-duality,sec:sva-transform,sec:sva-cross}）建立 6D 几何工具——twist/wrench、对偶功率、6×6 变换、运动/力叉积；\cref{sec:sva-inertia,sec:sva-accel} 给出空间惯性与空间加速度；\cref{sec:sva-lie} 把这一切统一到 $\mathfrak{se}(3)$ 的伴随表示，与 \cref{ch:lie} 双向接线。这些工具是本章后半段 $O(n)$ 算法的全部素材。

\section{Plücker 坐标：twist 与 wrench}
\label{sec:sva-plucker}

\paragraph{运动向量（twist）。}
刚体的瞬时运动由两个量完全确定：角速度 $\boldsymbol\omega\in\mathbb R^3$（绕瞬时轴的旋转快慢）与某参考点 $O$ 的线速度 $\mathbf v_O\in\mathbb R^3$（刚体上固连于 $O$ 之点的瞬时速度）。合并为 6D \emph{空间速度}（twist）

\begin{equation}
  \hat{\mathbf v}=\begin{pmatrix}\boldsymbol\omega\\ \mathbf v_O\end{pmatrix}\in\mathbb R^6,
  \qquad\text{（Featherstone 约定：角在前、线在后）}
  \label{eq:sva-twist}
\end{equation}
它生活在 6 维\emph{运动空间} $\mathcal M^6$。线性部分 $\mathbf v_O$ 依赖参考点选择：换到与 $O$ 相距 $\mathbf r_{OP}$ 的点 $P$，
\begin{equation}
  \mathbf v_P=\mathbf v_O+\boldsymbol\omega\times\mathbf r_{OP},
  \label{eq:sva-vref}
\end{equation}
而 $\boldsymbol\omega$ 与参考点无关。因此 twist 不是「自由向量」——其值取决于参考点，这正是 Plücker 坐标的本质特征。

\paragraph{与螺旋运动的关系。}
由 Mozzi--Chasles 定理，twist $\hat{\mathbf v}=(\boldsymbol\omega;\mathbf v_O)$ 编码了一个螺旋的全部信息：螺旋轴方向 $\boldsymbol\omega/\|\boldsymbol\omega\|$，离 $O$ 最近的轴上点 $\mathbf q=(\boldsymbol\omega\times\mathbf v_O)/\|\boldsymbol\omega\|^2$，角速度大小 $\|\boldsymbol\omega\|$，以及\emph{螺距} $h=(\boldsymbol\omega\cdot\mathbf v_O)/\|\boldsymbol\omega\|^2$（沿轴平移速度与旋转速度之比）。当 $\boldsymbol\omega=\mathbf 0$ 时退化为纯平移（螺距无穷）。twist 可在体坐标系（body-fixed）或空间坐标系（space-fixed）表达，二者由 \cref{sec:sva-transform} 的 6×6 变换相连；Featherstone 与多数库默认 body-fixed，因递推公式更简洁。

\paragraph{力向量（wrench）。}
对偶地，作用于刚体的外力由关于参考点 $O$ 的力矩 $\boldsymbol\tau_O\in\mathbb R^3$ 与合力 $\mathbf f\in\mathbb R^3$ 确定，合并为 6D \emph{力向量}（wrench）
\begin{equation}
  \hat{\mathbf f}=\begin{pmatrix}\boldsymbol\tau_O\\ \mathbf f\end{pmatrix}\in\mathbb R^6,
  \qquad\text{（力矩在前、力在后）}
  \label{eq:sva-wrench}
\end{equation}
它生活在另一 6 维\emph{力空间} $\mathcal F^6$。换参考点 $O\to P$ 时合力 $\mathbf f$ 不变，力矩变换为
\begin{equation}
  \boldsymbol\tau_P=\boldsymbol\tau_O-\mathbf r_{OP}\times\mathbf f.
  \label{eq:sva-tauref}
\end{equation}
推导：设 $\mathbf f$ 作用于点 $A$，则 $\boldsymbol\tau_O=\mathbf r_{OA}\times\mathbf f$，$\boldsymbol\tau_P=\mathbf r_{PA}\times\mathbf f=(\mathbf r_{PO}+\mathbf r_{OA})\times\mathbf f=\mathbf r_{PO}\times\mathbf f+\boldsymbol\tau_O$，代 $\mathbf r_{PO}=-\mathbf r_{OP}$ 即得。对比 \cref{eq:sva-vref}：twist 换参考点用 $+\boldsymbol\omega\times\mathbf r_{OP}$，wrench 换参考点用 $-\mathbf r_{OP}\times\mathbf f$，叉积方向相反——这正是对偶性的初次显现。

\begin{note}
\textbf{排序对照（[omega;v] vs 本书 $\boldsymbol\xi=[\boldsymbol\rho;\boldsymbol\phi]$）.}\;
本章 spatial vector 采 Featherstone $[\boldsymbol\omega;\mathbf v]$（角在前，算法文献主流），而 \cref{ch:lie} 的 $\mathfrak{se}(3)$ 坐标用 $\boldsymbol\xi=[\boldsymbol\rho;\boldsymbol\phi]$（线在前）。二者顺序相反，由置换矩阵
$\mathbf P=\begin{psmallmatrix}\mathbf 0&\mathbf I_3\\ \mathbf I_3&\mathbf 0\end{psmallmatrix}$
互换：$[\boldsymbol\omega;\mathbf v]=\mathbf P\,[\mathbf v;\boldsymbol\omega]$。注意 \cref{ch:lie} 还有 $\boldsymbol\rho\neq\mathbf t$（相差左雅可比）；本章的 $\mathbf v$ 是\emph{线速度}（运动学量），与 $\mathfrak{se}(3)$ 指数坐标的 $\boldsymbol\rho$ 不是同一对象，只有在「twist 作为李代数元素」时才需对齐排序。后文凡涉及 $\mathrm{ad}/\mathrm{ad}^*/\mathrm{Ad}$ 的块结构处再次明示此置换（\cref{sec:sva-lie}）。
\end{note}

\section{对偶性与功率标量积}
\label{sec:sva-duality}

\paragraph{功率的参考点无关性。}
twist 与 wrench 之间最根本的关系是：它们的标量积给出\emph{功率}，
\begin{equation}
  P=\hat{\mathbf f}^\top\hat{\mathbf v}=\boldsymbol\tau_O^\top\boldsymbol\omega+\mathbf f^\top\mathbf v_O,
  \label{eq:sva-power}
\end{equation}
而此值\textbf{不依赖参考点}。验证：换 $O\to P$，代入 \cref{eq:sva-tauref,eq:sva-vref}，
\[
  \boldsymbol\tau_P^\top\boldsymbol\omega+\mathbf f^\top\mathbf v_P
  =\boldsymbol\tau_O^\top\boldsymbol\omega-(\mathbf r_{OP}\times\mathbf f)^\top\boldsymbol\omega
   +\mathbf f^\top\mathbf v_O+\mathbf f^\top(\boldsymbol\omega\times\mathbf r_{OP}).
\]
用标量三重积 $(\mathbf a\times\mathbf b)\cdot\mathbf c=\mathbf a\cdot(\mathbf b\times\mathbf c)$ 与其轮换性：
\[
  -(\mathbf r_{OP}\times\mathbf f)\cdot\boldsymbol\omega=-\mathbf r_{OP}\cdot(\mathbf f\times\boldsymbol\omega)
  =\mathbf r_{OP}\cdot(\boldsymbol\omega\times\mathbf f),\qquad
  \mathbf f\cdot(\boldsymbol\omega\times\mathbf r_{OP})=\boldsymbol\omega\cdot(\mathbf r_{OP}\times\mathbf f)=-\mathbf r_{OP}\cdot(\boldsymbol\omega\times\mathbf f).
\]
两项一正一负相消，故 $P=\boldsymbol\tau_O^\top\boldsymbol\omega+\mathbf f^\top\mathbf v_O$ 与参考点无关。$\square$

\begin{insight}
运动空间 $\mathcal M^6$ 与力空间 $\mathcal F^6$ \textbf{不是同一向量空间}——它们是互相\emph{对偶}的空间，$\mathcal F^6=(\mathcal M^6)^*$。这不是学究式区分：对偶性决定了坐标变换规则\emph{不同}（运动向量用 $\mathbf X$，力向量用 $\mathbf X^{-\top}$，见 \cref{sec:sva-transform}）。若把 wrench 当 twist 来变换坐标，结果的物理含义就全错。功率 $P=\hat{\mathbf f}^\top\hat{\mathbf v}$ 是连接二者的唯一不变配对。
\end{insight}

\paragraph{对偶基。}
在坐标系 $\{O;\hat{\mathbf e}_1,\hat{\mathbf e}_2,\hat{\mathbf e}_3\}$ 下，$\mathcal M^6$ 有 6 个基向量：$\mathbf d_k=(\hat{\mathbf e}_k;\mathbf 0)$（$k=1,2,3$，沿三轴的单位旋转）与 $\mathbf d_{k+3}=(\mathbf 0;\hat{\mathbf e}_k)$（沿三轴的单位平移）。$\mathcal F^6$ 的对偶基 $\mathbf e^j$ 满足 $\langle\mathbf e^j,\mathbf d_k\rangle=\delta_{jk}$；在 Featherstone 约定下对偶基与 $\mathbf d_k$ 用相同的 6 元组表示。于是在同一坐标系下 twist 与 wrench 的分量是「相同的 6 元组」，却遵循不同变换规则——这与协变/逆变向量在正交基下分量相同、变换行为不同完全类比。

\section{6×6 Plücker 变换与力变换}
\label{sec:sva-transform}

\paragraph{运动变换的推导。}
设坐标系 $A$ 相对 $B$ 的姿态为 $\mathbf R\in\SO(3)$、位移为 $\mathbf p\in\mathbb R^3$（$A$ 原点相对 $B$ 原点，在 $B$ 系表达），即齐次变换 $\mathbf T_{BA}=\begin{psmallmatrix}\mathbf R&\mathbf p\\ \mathbf 0&1\end{psmallmatrix}$（建立在 \cref{ch:rigid_body} 的 SE(3) 之上）。给定 twist 在 $A$ 系的表达 ${}^A\hat{\mathbf v}=({}^A\boldsymbol\omega;{}^A\mathbf v_A)$，求其在 $B$ 系的表达。角速度只需旋转：${}^B\boldsymbol\omega=\mathbf R\,{}^A\boldsymbol\omega$。线速度需同时旋转并把参考点从 $A$ 原点换到 $B$ 原点（用 \cref{eq:sva-vref}，$\mathbf r_{OP}$ 即从 $A$ 原点到 $B$ 原点）：${}^B\mathbf v_B=\mathbf R\,{}^A\mathbf v_A+\mathbf p\times(\mathbf R\,{}^A\boldsymbol\omega)$。用 $\mathbf p\times(\mathbf R\boldsymbol\omega)=[\mathbf p]_\times\mathbf R\boldsymbol\omega$（$[\mathbf p]_\times$ 为反对称矩阵，见 \cref{ch:rigid_body}）合并为矩阵形式：

\begin{equation}
  {}^B\mathbf X_A=\begin{pmatrix}\mathbf R&\mathbf 0\\ [\mathbf p]_\times\mathbf R&\mathbf R\end{pmatrix},
  \qquad {}^B\hat{\mathbf v}={}^B\mathbf X_A\,{}^A\hat{\mathbf v}.
  \label{eq:sva-X}
\end{equation}
这就是 6×6 \emph{Plücker 变换矩阵}。其分块物理含义见 \cref{tab:sva-Xblocks}：右上零块表明\textbf{平移不影响角速度}（角速度是自由向量）。

\begin{table}[H]\centering\small
\caption{Plücker 运动变换 ${}^B\mathbf X_A$ 的分块结构与物理含义。}
\label{tab:sva-Xblocks}
\begin{tabular}{cll}
\toprule
分块 & 位置 & 含义 \\
\midrule
$\mathbf R$ & 左上 & 角速度分量的坐标系旋转 \\
$\mathbf 0$ & 右上 & 角速度不受平移影响（角速度是自由向量） \\
$[\mathbf p]_\times\mathbf R$ & 左下 & 参考点变换对线速度的耦合 \\
$\mathbf R$ & 右下 & 线速度分量的坐标系旋转 \\
\bottomrule
\end{tabular}
\end{table}

\paragraph{特例、逆与复合。}
纯旋转（$\mathbf p=\mathbf 0$）退化为对角分块 $\mathrm{diag}(\mathbf R,\mathbf R)$；纯平移（$\mathbf R=\mathbf I$）为下三角 $\begin{psmallmatrix}\mathbf I&\mathbf 0\\ [\mathbf p]_\times&\mathbf I\end{psmallmatrix}$，逆为 $\begin{psmallmatrix}\mathbf I&\mathbf 0\\ -[\mathbf p]_\times&\mathbf I\end{psmallmatrix}$。一般情况 $\mathbf X=\mathbf X_{\text{trans}}\mathbf X_{\text{rot}}$。逆变换把 $B$ 系表达变回 $A$ 系：
\begin{equation}
  {}^A\mathbf X_B=({}^B\mathbf X_A)^{-1}=\begin{pmatrix}\mathbf R^\top&\mathbf 0\\ -\mathbf R^\top[\mathbf p]_\times&\mathbf R^\top\end{pmatrix}.
  \label{eq:sva-Xinv}
\end{equation}
复合直接相乘 ${}^C\mathbf X_A={}^C\mathbf X_B\,{}^B\mathbf X_A$，与齐次变换 $\mathbf T_{CA}=\mathbf T_{CB}\mathbf T_{BA}$ 的复合规则一致——只是在 6D 空间中进行。

\paragraph{力变换 $\mathbf X^*=\mathbf X^{-\top}$。}
由功率不变性 \cref{eq:sva-power} 导出 wrench 的变换。设 ${}^A\hat{\mathbf v},{}^A\hat{\mathbf f}$ 是 $A$ 系中的运动/力对，要求 ${}^A\hat{\mathbf f}^\top{}^A\hat{\mathbf v}={}^B\hat{\mathbf f}^\top{}^B\hat{\mathbf v}$。代 ${}^B\hat{\mathbf v}={}^B\mathbf X_A\,{}^A\hat{\mathbf v}$，得 ${}^B\hat{\mathbf f}^\top{}^B\mathbf X_A={}^A\hat{\mathbf f}^\top$，对任意 ${}^A\hat{\mathbf v}$ 成立故
\begin{equation}
  {}^B\hat{\mathbf f}={}^B\mathbf X_A^{-\top}\,{}^A\hat{\mathbf f}=:\,{}^B\mathbf X_A^*\,{}^A\hat{\mathbf f},
  \qquad
  {}^B\mathbf X_A^*=\begin{pmatrix}\mathbf R&[\mathbf p]_\times\mathbf R\\ \mathbf 0&\mathbf R\end{pmatrix}.
  \label{eq:sva-Xstar}
\end{equation}
注意与运动变换的转置关系：力变换的右上角非零（力矩受平移影响），左下角为零（\textbf{合力是自由向量}，不受参考点变换影响）。

\begin{insight}
twist 与 wrench 的区别\textbf{不是}「一个有角量一个没有」——它们都有角量与线量。真正的区别是：twist 的角量（角速度）是自由向量、线量依赖参考点；wrench 恰好相反——合力是自由向量、力矩依赖参考点。这就是为什么变换矩阵的零块位置在 $\mathbf X$ 与 $\mathbf X^*$ 中互换（左下零 $\leftrightarrow$ 右上零）。
\end{insight}

\section{运动叉积与力叉积}
\label{sec:sva-cross}

\paragraph{运动叉积 $\times$。}
在运动坐标系中对空间向量求时间导数会产生一个修正项，称\emph{运动叉积}。设 $\hat{\mathbf v}_1=(\boldsymbol\omega_1;\mathbf v_1)$、$\hat{\mathbf v}_2=(\boldsymbol\omega_2;\mathbf v_2)$，定义
\begin{equation}
  \hat{\mathbf v}_1\times\hat{\mathbf v}_2:=\begin{pmatrix}\boldsymbol\omega_1\times\boldsymbol\omega_2\\ \boldsymbol\omega_1\times\mathbf v_2-\boldsymbol\omega_2\times\mathbf v_1\end{pmatrix}
  =[\hat{\mathbf v}_1]_\times\hat{\mathbf v}_2,
  \qquad
  [\hat{\mathbf v}]_\times=\begin{pmatrix}[\boldsymbol\omega]_\times&\mathbf 0\\ [\mathbf v]_\times&[\boldsymbol\omega]_\times\end{pmatrix}.
  \label{eq:sva-crm}
\end{equation}
$[\hat{\mathbf v}]_\times$ 是\emph{下三角}分块。其物理来源：若坐标系以 twist $\hat{\mathbf v}_1$ 运动，则在该系表达的空间向量 $\hat{\mathbf v}_2$ 的绝对时间导数为
$\frac{\mathrm d\hat{\mathbf v}_2}{\mathrm dt}\big|_{\text{fixed}}=\frac{\mathrm d\hat{\mathbf v}_2}{\mathrm dt}\big|_{\text{body}}+\hat{\mathbf v}_1\times\hat{\mathbf v}_2$，与三维公式 $\dot{\mathbf a}|_{\text{fixed}}=\dot{\mathbf a}|_{\text{rot}}+\boldsymbol\omega\times\mathbf a$ 完全类比，只是从 3D 提升到 6D。运动叉积满足：（i）双线性；（ii）反对称 $\hat{\mathbf v}_1\times\hat{\mathbf v}_2=-\hat{\mathbf v}_2\times\hat{\mathbf v}_1$；（iii）Jacobi 恒等式 $\hat{\mathbf v}_1\times(\hat{\mathbf v}_2\times\hat{\mathbf v}_3)+\text{(轮换)}=\mathbf 0$。这三条公理正是 $\mathfrak{se}(3)$ 李代数的结构（\cref{sec:sva-lie}）。

\paragraph{力叉积 $\times^*$。}
力叉积出现在动量定理的空间形式中，其定义由功率配对的不变性决定：对功率 $\hat{\mathbf f}^\top\hat{\mathbf v}$ 沿 $\hat{\mathbf v}_1$ 求时间导数并用 Leibniz 律，得余伴随关系 $\hat{\mathbf v}_2^\top(\hat{\mathbf v}_1\times^*\hat{\mathbf f})=-(\hat{\mathbf v}_1\times\hat{\mathbf v}_2)^\top\hat{\mathbf f}$（即 $\langle\hat{\mathbf v}_1\times^*\hat{\mathbf f},\hat{\mathbf v}_2\rangle=-\langle\hat{\mathbf f},\hat{\mathbf v}_1\times\hat{\mathbf v}_2\rangle$，与 \cref{thm:sva-lie} 的 $\mathrm{ad}^*=-\mathrm{ad}^\top$ 一致\footnote{此处负号不可省：$\mathfrak{se}(3)$ 的余伴随是 $\mathrm{ad}^*_{\hat{\mathbf v}}=-\mathrm{ad}_{\hat{\mathbf v}}^\top$（Featherstone\cite{featherstone2008rigid} \S2.9、Carpentier--Mansard\cite{carpentier2018analytical} 式 (2)–(3)）。若误把配对写成 $+(\hat{\mathbf v}_1\times\hat{\mathbf v}_2)^\top\hat{\mathbf f}$ 会得到符号相反的 $+[\hat{\mathbf v}]_\times^\top$，与下式及单刚体 Newton--Euler 的陀螺项 $\boldsymbol\omega\times(\mathbf I\boldsymbol\omega)$ 矛盾。}），即 $[\hat{\mathbf v}]_{\times^*}=-[\hat{\mathbf v}]_\times^\top$，故
\begin{equation}
  [\hat{\mathbf v}]_{\times^*}=-[\hat{\mathbf v}]_\times^\top=\begin{pmatrix}[\boldsymbol\omega]_\times&[\mathbf v]_\times\\ \mathbf 0&[\boldsymbol\omega]_\times\end{pmatrix}.
  \label{eq:sva-crf}
\end{equation}
\textbf{运动叉积是下三角，力叉积是上三角}——又一处对偶的零块互换。展开 $\hat{\mathbf v}\times^*\hat{\mathbf f}=(\boldsymbol\omega\times\boldsymbol\tau+\mathbf v\times\mathbf f;\;\boldsymbol\omega\times\mathbf f)$：力矩变化 = 角速度致力矩旋转 + 平移速度与力的耦合，合力变化 = 角速度致力方向旋转。\cref{tab:sva-cross-explicit} 给出两个 6×6 矩阵的完整元素（经联网与 Carpentier--Mansard\cite{carpentier2018analytical} 逐元素核对）。

\begin{table}[H]\centering\footnotesize
\caption{运动叉积矩阵 $[\hat{\mathbf v}]_\times$（下三角）与力叉积矩阵 $[\hat{\mathbf v}]_{\times^*}=-[\hat{\mathbf v}]_\times^\top$（上三角），$\hat{\mathbf v}=(\omega_x,\omega_y,\omega_z,v_x,v_y,v_z)$。}
\label{tab:sva-cross-explicit}
\renewcommand{\arraystretch}{1.15}
\begin{tabular}{cc}
$[\hat{\mathbf v}]_\times=\begin{psmallmatrix}
0&-\omega_z&\omega_y&0&0&0\\
\omega_z&0&-\omega_x&0&0&0\\
-\omega_y&\omega_x&0&0&0&0\\
0&-v_z&v_y&0&-\omega_z&\omega_y\\
v_z&0&-v_x&\omega_z&0&-\omega_x\\
-v_y&v_x&0&-\omega_y&\omega_x&0
\end{psmallmatrix}$
&
$[\hat{\mathbf v}]_{\times^*}=\begin{psmallmatrix}
0&-\omega_z&\omega_y&0&-v_z&v_y\\
\omega_z&0&-\omega_x&v_z&0&-v_x\\
-\omega_y&\omega_x&0&-v_y&v_x&0\\
0&0&0&0&-\omega_z&\omega_y\\
0&0&0&\omega_z&0&-\omega_x\\
0&0&0&-\omega_y&\omega_x&0
\end{psmallmatrix}$
\end{tabular}
\end{table}

\section{空间惯性张量}
\label{sec:sva-inertia}

\paragraph{空间动量与质心系构造。}
在空间框架中，存在 6×6 矩阵 $\mathcal I$ 使空间动量 $\hat{\mathbf h}=\mathcal I\hat{\mathbf v}$（$\hat{\mathbf h}$ 是 wrench：角动量在前、线动量在后）。若取质心为参考点，
\begin{equation}
  \mathcal I_c=\begin{pmatrix}\mathbf I_c&\mathbf 0\\ \mathbf 0&m\mathbf I_3\end{pmatrix},
  \label{eq:sva-Ic}
\end{equation}
对角分块——角动量部分为 $\mathbf I_c\boldsymbol\omega$、线动量部分为 $m\mathbf v_c$，二者解耦。

\paragraph{非质心参考点（平行轴空间版）。}
若参考点 $O$ 偏离质心，记 $\mathbf c$ 为从 $O$ 到质心的位置矢量。设 twist 在 $O$ 表达为 $(\boldsymbol\omega;\mathbf v_O)$，则质心速度 $\mathbf v_c=\mathbf v_O+\boldsymbol\omega\times\mathbf c$。线动量
$\mathbf p=m\mathbf v_c=m\mathbf v_O+m(\boldsymbol\omega\times\mathbf c)=m\mathbf v_O-m[\mathbf c]_\times\boldsymbol\omega$；角动量
\[
  \mathbf L_O=\mathbf I_c\boldsymbol\omega+\mathbf c\times(m\mathbf v_c)
  =\mathbf I_c\boldsymbol\omega+m[\mathbf c]_\times\mathbf v_O+m[\mathbf c]_\times(\boldsymbol\omega\times\mathbf c)
  =(\mathbf I_c-m[\mathbf c]_\times[\mathbf c]_\times)\boldsymbol\omega+m[\mathbf c]_\times\mathbf v_O.
\]
合并 $\hat{\mathbf h}_O=(\mathbf L_O;\mathbf p)$ 即得（与 Featherstone RBDA (2.63) 一致\cite{featherstone2008rigid}）：
\begin{equation}
  \mathcal I_O=\begin{pmatrix}\mathbf I_c-m[\mathbf c]_\times[\mathbf c]_\times & m[\mathbf c]_\times\\ -m[\mathbf c]_\times & m\mathbf I_3\end{pmatrix}
  =\begin{pmatrix}\bar{\mathbf I} & m[\mathbf c]_\times\\ m[\mathbf c]_\times^\top & m\mathbf I_3\end{pmatrix},
  \label{eq:sva-IO}
\end{equation}
其中 $\bar{\mathbf I}=\mathbf I_c+m(\mathbf c^\top\mathbf c\,\mathbf I_3-\mathbf c\mathbf c^\top)$ 是平行轴修正后的转动惯量（用了 $-[\mathbf c]_\times[\mathbf c]_\times=\|\mathbf c\|^2\mathbf I-\mathbf c\mathbf c^\top$）。右上 $m[\mathbf c]_\times$ 与左下 $m[\mathbf c]_\times^\top=-m[\mathbf c]_\times$ 互为转置——\textbf{$\mathcal I$ 是对称矩阵}，这不是巧合，而是动能 $T=\tfrac12\hat{\mathbf v}^\top\mathcal I\hat{\mathbf v}$ 必为二次型的直接要求。$\mathcal I$ 由 10 个独立参数刻画：质量 $m$（1）、质心 $\mathbf c$（3）、$\mathbf I_c$（6）。

\paragraph{平行轴定理的算子形式。}
由动能在坐标变换下不变 $T=\tfrac12(\mathbf X\hat{\mathbf v}_A)^\top\mathcal I_B(\mathbf X\hat{\mathbf v}_A)=\tfrac12\hat{\mathbf v}_A^\top\mathcal I_A\hat{\mathbf v}_A$，要求 $\mathbf X^\top\mathcal I_B\mathbf X=\mathcal I_A$，即
\begin{equation}
  \boxed{\;\mathcal I_B=\mathbf X^{-\top}\mathcal I_A\mathbf X^{-1}\;}
  \qquad(\mathbf X={}^B\mathbf X_A).
  \label{eq:sva-Iparallel}
\end{equation}
空间惯性像张量一样被逆变换的转置从两侧「夹」。这一规则与 \cref{eq:sva-Xstar} 的力变换 $\mathbf X^*=\mathbf X^{-\top}$ 一致（$\mathcal I$ 把 twist 映为 wrench，故两侧各用一次力/运动变换）。

\paragraph{物理一致性。}
$\mathcal I$ 对称是动能二次型的必要条件，但不充分。\emph{完全物理一致性}（full physical consistency，Wensing 等\cite{wensing2018linear}）还要求修正后的转动惯量 $\bar{\mathbf I}$ 的三个主惯量 $I_1,I_2,I_3$ 满足\textbf{三角不等式} $I_i+I_j\ge I_k$（对所有排列）——这等价于存在一个真实质量分布产生该惯性。仅检查正定性会漏掉违反三角不等式的非物理参数（如 URDF 标注错误），导致仿真不稳定。复合刚体惯性的累加规则（为 \cref{sec:on-crba} 的 CRBA 铺垫）极其简单：把子连杆惯性用 \cref{eq:sva-Iparallel} 变换到当前坐标系后直接相加，
\begin{equation}
  \mathcal I_i^c=\mathcal I_i+\sum_{j\in\text{children}(i)}{}^i\mathbf X_j^*\,\mathcal I_j^c\,{}^j\mathbf X_i.
  \label{eq:sva-Icomposite}
\end{equation}

\section{空间加速度与单刚体 Newton--Euler 统一}
\label{sec:sva-accel}

\paragraph{空间加速度 vs 经典加速度。}
空间加速度定义为空间速度的时间导数（在固定系中）$\hat{\mathbf a}=\mathrm d\hat{\mathbf v}/\mathrm dt=(\dot{\boldsymbol\omega};\dot{\mathbf v}_O)$；在体系表达需加运动叉积修正 $\hat{\mathbf a}_{\text{body}}=\frac{\mathrm d\hat{\mathbf v}}{\mathrm dt}|_{\text{body}}+\hat{\mathbf v}_{\text{frame}}\times\hat{\mathbf v}$。它与传感器直接测得的\emph{经典加速度}的关系是
\begin{equation}
  \hat{\mathbf a}_{\text{cls}}=\hat{\mathbf a}_{\text{sp}}+\begin{pmatrix}\mathbf 0\\ \boldsymbol\omega\times\mathbf v\end{pmatrix},
  \label{eq:sva-accel-cls}
\end{equation}
角加速度分量相同，差别只在线性部分 $\boldsymbol\omega\times\mathbf v$。原因：$\dot{\mathbf v}_O$ 是「瞬时固定在空间中的参考点 $O$」的加速度在体系的表达，而经典加速度是「刚体上物理附着于 $O$ 的质点」的加速度——下一瞬间体系已旋转，「$O$ 点」不再是同一质点。

\begin{insight}
\textbf{为何空间加速度无 Coriolis 项？}（这是空间向量代数最不直觉但最重要的特性。）Featherstone 的总结是：「空间加速度是空间速度的时间导数，因此服从与速度相同的加法规则……没有 Coriolis 项，因为空间加速度（不同于经典加速度）是一个\emph{真正的向量}」——它在坐标变换下像向量一样变换，不需额外修正项。这与流体力学的 Euler 描述（在空间固定点观察流过的流体）对应，经典加速度对应 Lagrange 描述（跟踪特定流体微元），二者之差正是对流项 $\boldsymbol\omega\times\mathbf v$。这一性质使后文 RNEA 的加速度递推只需一个运动叉积项、不需展开 Coriolis 与离心加速度（\cref{sec:on-rnea}）。
\end{insight}

\paragraph{单刚体空间 Newton--Euler 方程。}
经典 Newton--Euler 在三维下是两个独立方程 $\mathbf f=m\mathbf a_c$、$\boldsymbol\tau_c=\mathbf I_c\dot{\boldsymbol\omega}+\boldsymbol\omega\times(\mathbf I_c\boldsymbol\omega)$。在质心系工作（$\mathcal I_c=\mathrm{diag}(\mathbf I_c,m\mathbf I_3)$、$\hat{\mathbf v}=(\boldsymbol\omega;\mathbf v_c)$、$\hat{\mathbf a}=(\dot{\boldsymbol\omega};\mathbf a_c)$），二者统一为一式：
\begin{equation}
  \boxed{\;\hat{\mathbf f}=\mathcal I\hat{\mathbf a}+\hat{\mathbf v}\times^*(\mathcal I\hat{\mathbf v})\;}
  \label{eq:sva-NE}
\end{equation}
验证力叉积项：
$\hat{\mathbf v}\times^*(\mathcal I_c\hat{\mathbf v})=\big(\boldsymbol\omega\times(\mathbf I_c\boldsymbol\omega)+\mathbf v_c\times(m\mathbf v_c);\;\boldsymbol\omega\times(m\mathbf v_c)\big)=\big(\boldsymbol\omega\times(\mathbf I_c\boldsymbol\omega);\;m\boldsymbol\omega\times\mathbf v_c\big)$（$\mathbf v_c\times\mathbf v_c=\mathbf 0$）。加惯性项 $\mathcal I_c\hat{\mathbf a}=(\mathbf I_c\dot{\boldsymbol\omega};m\mathbf a_c)$ 得 $\hat{\mathbf f}=(\mathbf I_c\dot{\boldsymbol\omega}+\boldsymbol\omega\times(\mathbf I_c\boldsymbol\omega);\;m(\mathbf a_c+\boldsymbol\omega\times\mathbf v_c))$。力矩分量正是 Euler 方程；力分量 $m(\mathbf a_c+\boldsymbol\omega\times\mathbf v_c)=m\mathbf a_{\text{cls}}=\mathbf f$（用 \cref{eq:sva-accel-cls}），正是 Newton 第二定律。一个方程统一了两个。若参考点不在质心，$\mathcal I_O$ 有非对角块，但方程形式完全不变——所有「平行轴修正」被吸收进 $\mathcal I_O$。

\begin{insight}
\cref{eq:sva-NE} 中的力叉积项 $\hat{\mathbf v}\times^*(\mathcal I\hat{\mathbf v})$ 自动包含了三维方程的陀螺项 $\boldsymbol\omega\times(\mathbf I\boldsymbol\omega)$（离心/Coriolis 效应）与空间加速度到经典加速度的修正。换言之，\textbf{Coriolis 力没有消失，它隐藏在力叉积里}——这一观察是后文「为何 RNEA 不需显式构造 $\mathbf C$ 矩阵」与几何力学章 Christoffel 符号显式构造（\cref{ch:geometric-mechanics}）两视角的接口：本章给「Coriolis 在 $\hat{\mathbf v}\times^*\mathcal I\hat{\mathbf v}$ 中」的递推视角，几何力学章给 $C_{ijk}$ 的解析视角，同一对象两面。
\end{insight}

\section{李代数视角：ad、ad* 与 Ad 的统一}
\label{sec:sva-lie}

\paragraph{运动空间即 $\mathfrak{se}(3)$。}
前几节的全部规则，本质上是 $SE(3)$ 李群结构（\cref{ch:lie}）在 6D 刚体动力学上的具体化。从纯数学看，运动空间配运动叉积 $(\mathcal M^6,\times)$ 构成李代数，同构于 $\mathfrak{se}(3)$：李代数三公理（双线性、反对称、Jacobi）恰是 \cref{sec:sva-cross} 验证过的运动叉积性质（Jacobi 恒等式由 $\mathbb R^3$ 叉积的 Jacobi 恒等式直接展开得证）。

\begin{theorem}[空间向量算子与 $SE(3)$ 表示的对应]\label{thm:sva-lie}
设 $\mathbf T_{BA}\in SE(3)$，$\hat{\mathbf v}\in\mathcal M^6$，$\hat{\mathbf f}\in\mathcal F^6=\mathfrak{se}(3)^*$。则
\begin{enumerate}
  \item Plücker 运动变换是伴随表示：${}^B\mathbf X_A=\mathrm{Ad}_{\mathbf T_{BA}}$，且 ${}^C\mathbf X_A={}^C\mathbf X_B\,{}^B\mathbf X_A\iff\mathrm{Ad}_{\mathbf T_{CA}}=\mathrm{Ad}_{\mathbf T_{CB}}\!\circ\mathrm{Ad}_{\mathbf T_{BA}}$（$\mathrm{Ad}$ 是群同态）；
  \item 力变换是余伴随表示：${}^B\mathbf X_A^*=\mathbf X^{-\top}=\mathrm{Ad}^*_{\mathbf T_{BA}}$，由 $\langle\mathrm{Ad}^*_g\boldsymbol\mu,\boldsymbol\xi\rangle=\langle\boldsymbol\mu,\mathrm{Ad}_{g^{-1}}\boldsymbol\xi\rangle$ 定义；
  \item 运动叉积是伴随作用：$[\hat{\mathbf v}]_\times=\mathrm{ad}_{\hat{\mathbf v}}$；力叉积是余伴随作用：$[\hat{\mathbf v}]_{\times^*}=\mathrm{ad}^*_{\hat{\mathbf v}}=-\mathrm{ad}_{\hat{\mathbf v}}^\top$；
  \item 功率配对即李代数与其对偶的自然配对：$\langle\hat{\mathbf f},\hat{\mathbf v}\rangle=\hat{\mathbf f}^\top\hat{\mathbf v}=P$。
\end{enumerate}
\end{theorem}

\begin{proof}[证明骨架]
（1）伴随表示定义为群上共轭作用 $\mathrm{Ad}_g(\boldsymbol\xi)=g\boldsymbol\xi g^{-1}$ 降到李代数；把 $\boldsymbol\xi$ 写成 4×4 矩阵 $\begin{psmallmatrix}[\boldsymbol\omega]_\times&\mathbf v\\ \mathbf 0&0\end{psmallmatrix}$，对 $\mathbf T=\begin{psmallmatrix}\mathbf R&\mathbf p\\ \mathbf 0&1\end{psmallmatrix}$ 直接计算 $\mathbf T\boldsymbol\xi\mathbf T^{-1}$，其 $(\boldsymbol\omega',\mathbf v')$ 系数恰为 \cref{eq:sva-X}。群同态性由共轭复合 $g_1g_2(\cdot)(g_1g_2)^{-1}$ 给出。（2）余伴随由对偶配对的不变性定义，转置关系 $\mathrm{Ad}^*=\mathrm{Ad}^{-\top}$ 即 \cref{eq:sva-Xstar}。（3）$\mathrm{ad}_{\hat{\mathbf v}_1}\hat{\mathbf v}_2=[\hat{\mathbf v}_1,\hat{\mathbf v}_2]$ 即李括号=运动叉积，矩阵形式 \cref{eq:sva-crm}；余伴随 $\mathrm{ad}^*=-\mathrm{ad}^\top$ 由 $\langle\mathrm{ad}^*_{\hat{\mathbf v}}\hat{\mathbf f},\hat{\mathbf u}\rangle=-\langle\hat{\mathbf f},\mathrm{ad}_{\hat{\mathbf v}}\hat{\mathbf u}\rangle$ 与 \cref{eq:sva-crf} 一致。（4）显然。
\end{proof}

\cref{tab:sva-lie-dict} 汇总这一对应。它的实践价值：一旦理解对偶与伴随，所有变换规则都从一个统一原理（群作用在自身与对偶空间上的表示）自动导出——无需分别记忆运动变换与力变换，只需记「运动用 $\mathrm{Ad}$，力用 $\mathrm{Ad}^*$，功率配对是不变量」。指数/对数映射、右扰动 $\mathbf J_r$ 的来源在 \cref{ch:lie}；本章把抽象的 $\mathrm{ad}/\mathrm{ad}^*$ 具体化为 6D 刚体动力学，故 \cref{ch:lie} 可反向交叉引用本章作为「伴随表示的动力学应用」。

\begin{table}[H]\centering\small
\caption{空间向量算子 $\leftrightarrow$ $\mathfrak{se}(3)$ 表示的对应字典。}
\label{tab:sva-lie-dict}
\begin{tabular}{lll}
\toprule
数学概念 & 空间向量记号 & 作用对象 \\
\midrule
$\mathfrak{se}(3)$ & $\mathcal M^6$（运动空间） & twist \\
$\mathfrak{se}(3)^*$ & $\mathcal F^6$（力空间） & wrench \\
$\mathrm{Ad}_{\mathbf T}$ & $\mathbf X$（Plücker 变换） & 运动坐标变换 \\
$\mathrm{Ad}^*_{\mathbf T}$ & $\mathbf X^{-\top}=\mathbf X^*$ & 力坐标变换 \\
$\mathrm{ad}_{\hat{\mathbf v}}$ & $[\hat{\mathbf v}]_\times$（运动叉积，下三角） & twist \\
$\mathrm{ad}^*_{\hat{\mathbf v}}$ & $[\hat{\mathbf v}]_{\times^*}=-[\hat{\mathbf v}]_\times^\top$（力叉积，上三角） & wrench \\
配对 $\langle\boldsymbol\mu,\boldsymbol\xi\rangle$ & $\hat{\mathbf f}^\top\hat{\mathbf v}$ & 功率 \\
\bottomrule
\end{tabular}
\end{table}

\begin{note}
\textbf{排序冲突在此再次明示.}\;
\cref{thm:sva-lie} 的块结构\emph{依赖排序约定}。本章 $[\boldsymbol\omega;\mathbf v]$ 下 $\mathrm{ad}_{\hat{\mathbf v}}=\begin{psmallmatrix}[\boldsymbol\omega]_\times&\mathbf 0\\ [\mathbf v]_\times&[\boldsymbol\omega]_\times\end{psmallmatrix}$（下三角）；而本书 \cref{ch:lie} 用 $\boldsymbol\xi=[\boldsymbol\rho;\boldsymbol\phi]$（线在前），同一算子写作 $\mathrm{ad}_{\boldsymbol\xi}=\begin{psmallmatrix}[\boldsymbol\phi]_\times&[\boldsymbol\rho]_\times\\ \mathbf 0&[\boldsymbol\phi]_\times\end{psmallmatrix}$——\textbf{块的上下三角位置因排序不同而互换}。二者经 \cref{sec:sva-plucker} 的置换矩阵 $\mathbf P$ 相联：$\mathrm{ad}^{[\omega;v]}=\mathbf P\,\mathrm{ad}^{[\rho;\phi]}\,\mathbf P$。跨章引用 $\mathrm{ad}$ 块结构时务必先核对排序，这是本章最易错处。
\end{note}

\paragraph{三流派记号对照。}
读不同文献需在三套记号间互译，见 \cref{tab:sva-notations}。Featherstone 记号算法导向、被多数库内部采用；Lynch--Park\cite{lynch2017modern,murray1994mathematical} 强调 $SE(3)$ 李群结构、与微分几何/控制语言统一；Drake 用 monogram 命名 \texttt{V\_MB\_E}（B 相对 M、在 E 中表达），且明确声明其空间向量「不是 Plücker 向量」、用显式 \texttt{Shift}/\allowbreak\texttt{Rotate} 分步操作而非 6×6 矩阵乘——这是跨库迁移代码时最易踩的坑（直接套用 6×6 公式会得错误结果）。

\begin{table}[H]\centering\small
\caption{三流派空间向量记号对照（读论文/跨库互译工具）。}
\label{tab:sva-notations}
\begin{tabular}{llll}
\toprule
概念 & Featherstone & Lynch--Park & Drake \\
\midrule
体速度 & $\hat{\mathbf v}=(\boldsymbol\omega;\mathbf v)$ & $\mathcal V_b=(\boldsymbol\omega_b;\mathbf v_b)$ & \texttt{V\_WB\_B} \\
坐标变换 & ${}^B\mathbf X_A\hat{\mathbf v}$ & $[\mathrm{Ad}_{T_{BA}}]\mathcal V$ & \texttt{V.\allowbreak Shift(p).\allowbreak Rotate(R)} \\
力变换 & $\mathbf X^*\hat{\mathbf f}=\mathbf X^{-\top}\hat{\mathbf f}$ & $[\mathrm{Ad}_T]^{-\top}\mathcal F$ & \texttt{F.Shift(p)} \\
运动叉积 & $\hat{\mathbf v}_1\times\hat{\mathbf v}_2$ & $[\mathrm{ad}_{\mathcal V_1}]\mathcal V_2$ & 不显式暴露 \\
惯性变换 & $\mathbf X^{-\top}\mathcal I\mathbf X^{-1}$ & $[\mathrm{Ad}_T]^\top\mathcal G[\mathrm{Ad}_T]$ & 内部处理 \\
Newton--Euler & $\hat{\mathbf f}=\mathcal I\hat{\mathbf a}+\hat{\mathbf v}\times^*\mathcal I\hat{\mathbf v}$ & $\mathcal F=\mathcal G\dot{\mathcal V}-[\mathrm{ad}_{\mathcal V}]^\top\mathcal G\mathcal V$ & \texttt{CalcInverseDynamics} \\
\bottomrule
\end{tabular}
\end{table}

\paragraph{承上启下。}
至此 6D 几何工具齐备：twist/wrench（\cref{sec:sva-plucker}）、对偶功率（\cref{sec:sva-duality}）、6×6 变换 $\mathbf X/\mathbf X^*$（\cref{sec:sva-transform}）、运动/力叉积 $\times/\times^*$（\cref{sec:sva-cross}）、空间惯性 $\mathcal I$（\cref{sec:sva-inertia}）、空间加速度与单体 Newton--Euler（\cref{sec:sva-accel}）、以及它们到 $\mathrm{ad}/\mathrm{ad}^*/\mathrm{Ad}$ 的统一（本节）。本章后半段将看到，这些工具恰好是 $O(n)$ 递推算法的全部素材——RNEA 的速度递推用 $\mathbf X$ 与 $\times$，力回溯用 $\times^*$ 与 $\mathbf X^*$，CRBA 用 $\mathcal I^c$ 累加，ABA 用 $\mathcal I$ 的 Schur 补。下一节从算法史进入。

\section{算法史与复杂度演进}
\label{sec:on-history}

\paragraph{动机：实时的刚性需求。}
2008 年至今，几乎所有主流机器人软件栈的动力学内核都是 Featherstone 三大算法 RNEA/CRBA/ABA 的某种实现——Pinocchio\cite{carpentier2019pinocchio}、RBDL、Drake、MuJoCo、iDynTree 等库的内部循环逐行映射到 RBDA\cite{featherstone2008rigid} 第 5、6、7 章的三张伪代码。其重要性源于实时控制的刚性约束：现代优化控制（如 \cref{ch:ddp-ilqr} 的 DDP/iLQR、\cref{ch:mpc-advanced} 的 MPC）每次迭代约调用 $N$ 次 ABA $+\,N$ 次 RNEA $+\,N$ 次解析导数（$N$ 是 horizon 节点数，典型 60--115）。在 36 自由度机器人上，若动力学每次调用 $>10\,\mu\mathrm s$，单次 MPC 迭代就超过 10\,ms，无法满足 100\,Hz 以上闭环。微秒级动力学不是优化，而是可行性的前提。

\paragraph{从 $O(n^4)$ 到 $O(n)$。}
早期 Uicker--Kahn 方法（1965）按定义逐元素算 $\mathbf M(\mathbf q)$ 与 Christoffel 符号，复杂度 $O(n^4)$，完全不利用关节树的递推结构；30 自由度机器人需约 81 万次浮点运算。此后四十年是「三个数量级加速」的历史（\cref{tab:on-history}，Featherstone--Orin ICRA 2000 综述\cite{featherstone2000robot}）。值得澄清一处常见史料错误：Hollerbach 的递推 Lagrange 论文实为 \textbf{1980} 年（非 1978），且其核心贡献恰是把 Lagrange 公式的运算量从旧式 $O(n^4)$ 降到\textbf{随关节数线性增长}（即 $O(n)$，只是常数较大），并非 $O(n^3)$。因此 $O(n^4)\to O(n)$ 的突破由两条\emph{同年}路线达成：Hollerbach（1980，递推 Lagrange 路线）\cite{hollerbach1980recursive} 与 Luh--Walker--Paul（1980，Newton--Euler 路线，即 RNEA）\cite{luh1980online}，二者殊途同归——Silver（1982）证明两种公式无本质效率差异\cite{silver1982equivalence}。

\begin{table}[H]\centering\small
\caption{动力学算法复杂度演进（已订正 Hollerbach 1980/$O(n)$）。}
\label{tab:on-history}
\begin{tabular}{lllll}
\toprule
年代 & 方法 & 复杂度 & 作者 & 核心思想 \\
\midrule
1965 & Uicker--Kahn & $O(n^4)$ & Uicker & 逐元素算 $\mathbf M$ 与 Christoffel 符号 \\
1980 & 递推 Lagrange & $O(n)$（常数大） & Hollerbach & 首次把 Lagrange 降到线性 \\
1980 & RNEA & $O(n)$ & Luh, Walker, Paul & Newton--Euler 前向+后向递推 \\
1982 & CRBA & $O(n^2)$ & Walker, Orin & 复合刚体惯性后向累加 \\
1983 & ABA & $O(n)$ & Featherstone & 关节体惯性 + Schur 消元 \\
\bottomrule
\end{tabular}
\end{table}

对 30 自由度人形：$O(n^4)$ 需约 $810{,}000$ FLOPs，$O(n)$ 只需约 $6{,}600$ FLOPs，相差约 123 倍。

\begin{insight}
递推算法能达到 $O(n)$ 的根本原因是\textbf{树状拓扑的局部性}——每个关节只需与父/子关节通信，全局问题分解为 $n$ 个局部子问题的串联（类比卷积是局部操作故可并行）。一旦机器人含闭链（如平行四边形机构），局部性被打破，算法必须引入约束力 $\boldsymbol\lambda$，回到 $O(n)+O(m^3)$ 形式（$m$ 为约束数）——这正是约束动力学章（\cref{ch:constrained-dynamics}）的起点。
\end{insight}

\section{关节树拓扑与运动子空间}
\label{sec:on-tree}

机器人模型在数学上是一棵\emph{有向连通树}：节点是刚体（link），边是关节（joint）。三个关键定义：

\begin{itemize}
  \item \textbf{父数组 $\lambda(i)$}：关节 $i$ 的父关节编号，基座取 $\lambda=0$；
  \item \textbf{拓扑序编号}：关节从 1 到 $n$ 编号，\textbf{保证 $\lambda(i)<i$}（父节点编号小于子节点）。这使「从根到叶」遍历可用 \texttt{for $i=1..\allowbreak n$}、「从叶到根」可用 \texttt{for $i=n..1$}；
  \item \textbf{运动子空间 $\mathbf S_i\in\mathbb R^{6\times n_i}$}（$n_i$ 为关节 $i$ 的自由度数）：描述关节允许的运动方向，见 \cref{tab:on-S}。
\end{itemize}

拓扑序的必要性是\emph{因果性}：前向递推 $\mathbf v_i={}^i\mathbf X_{\lambda(i)}\mathbf v_{\lambda(i)}+\mathbf S_i\dot q_i$ 用到父节点速度，若 $\lambda(i)>i$ 则计算 $\mathbf v_i$ 时父速度尚未算出，数据依赖被打破。这与计算图的拓扑排序完全一致——深度学习框架的前向传播也必须按拓扑序执行有向无环图的节点。$\mathbf S_i$ 与全局几何 Jacobian $\mathbf J$ 的关系为 $\mathbf J=[\cdots\,{}^0\mathbf X_i\mathbf S_i\,\cdots]$，机械臂部 Jacobian 章在此对齐。

\begin{table}[H]\centering\small
\caption{各关节类型的运动子空间 $\mathbf S_i$（z 轴约定）。}
\label{tab:on-S}
\begin{tabular}{lcl}
\toprule
关节类型 & $n_i$ & $\mathbf S_i$ \\
\midrule
Revolute（绕 z） & 1 & $(0,0,1,0,0,0)^\top$ \\
Prismatic（沿 z） & 1 & $(0,0,0,0,0,1)^\top$ \\
Spherical & 3 & $\begin{psmallmatrix}\mathbf I_3\\ \mathbf 0_3\end{psmallmatrix}$ \\
Free/Float（浮基） & 6 & $\mathbf I_6$ \\
\bottomrule
\end{tabular}
\end{table}

\section{RNEA：逆动力学 \texorpdfstring{$O(n)$}{O(n)}}
\label{sec:on-rnea}

\paragraph{问题与动机。}
RNEA（递推 Newton--Euler 算法）求解逆动力学：已知 $\mathbf q,\dot{\mathbf q},\ddot{\mathbf q}$，求所需关节力矩
\begin{equation}
  \boldsymbol\tau=\mathrm{RNEA}(\mathbf q,\dot{\mathbf q},\ddot{\mathbf q})=\mathbf M(\mathbf q)\ddot{\mathbf q}+\mathbf C(\mathbf q,\dot{\mathbf q})\dot{\mathbf q}+\mathbf g(\mathbf q).
  \label{eq:on-rnea-def}
\end{equation}
若直接用 Lagrange 公式，构造 $\mathbf M$ 本身就是 $O(n^2)$（CRBA），再做 $\mathbf M\ddot{\mathbf q}$ 又 $O(n^2)$；而 RNEA 根本不显式构造 $\mathbf M,\mathbf C,\mathbf g$，直接从 Newton--Euler 方程出发、经树上递推一次得到 $\boldsymbol\tau$，复杂度 $O(n)$。

\paragraph{前向 pass：速度与加速度递推。}
从根向叶传播。空间速度递推
\begin{equation}
  \mathbf v_i={}^i\mathbf X_{\lambda(i)}\mathbf v_{\lambda(i)}+\mathbf S_i\dot q_i,
  \label{eq:on-vrec}
\end{equation}
含义为「父速度变换到本系 + 关节自身贡献」。空间加速度递推
\begin{equation}
  \mathbf a_i={}^i\mathbf X_{\lambda(i)}\mathbf a_{\lambda(i)}+\mathbf S_i\ddot q_i+\underbrace{\mathbf v_i\times(\mathbf S_i\dot q_i)}_{\text{bias accel.}\;\mathbf c_i},
  \label{eq:on-arec}
\end{equation}
其中 $\mathbf c_i=\mathbf v_i\times(\mathbf S_i\dot q_i)$ 是\emph{偏置加速度}，本质是旋转坐标系中的科氏效应：即使 $\ddot q_i=0$，只要 link 在转（$\mathbf v_i\neq\mathbf 0$）且关节在动（$\dot q_i\neq0$），link 上的点就有加速度——与地球科氏力让风偏转是同一物理机制。

\paragraph{重力注入技巧。}
基座初始化 $\mathbf v_0=\mathbf 0$、$\mathbf a_0=-\mathbf g_0$，其中 $\mathbf g_0=(0,0,0,\mathbf g^\top)^\top$（重力的 spatial 表示，角部分零、线部分为重力加速度 $\mathbf g$，方向/符号统一于全书重力约定）。给基座一个「虚拟向上加速度」，等效于让所有 link 感受向下重力，从而\emph{不必}在力回溯中对每个 link 单独加 $m_i\mathbf g$。这把「分布在 $n$ 个 link 上的重力」压缩成「基座一次赋值」，极其优雅。须注意这是\emph{空间加速度}约定下的注入；若误用经典加速度约定，$\mathbf g$ 的注入会多一个交叉项（\cref{sec:sva-accel}）。

\paragraph{反向 pass：力回溯。}
从叶向根传播。单 link 的 Newton--Euler 力（即 \cref{eq:sva-NE}）
\begin{equation}
  \mathbf f_i=\mathcal I_i\mathbf a_i+\mathbf v_i\times^*(\mathcal I_i\mathbf v_i)-\mathbf f_i^{\text{ext}},
  \label{eq:on-frec}
\end{equation}
其中 $\mathbf v_i\times^*(\mathcal I_i\mathbf v_i)$ 是 velocity product force（含科氏力与离心力），$\mathbf f_i^{\text{ext}}$ 是外力（接触力可作为额外空间力在此注入，见 \cref{ch:contact-mechanics}）。力向父节点回溯（达朗贝尔原理）$\mathbf f_{\lambda(i)}\mathrel{+}={}^{\lambda(i)}\mathbf X_i^*\mathbf f_i$，关节力矩为净力在运动方向的投影 $\tau_i=\mathbf S_i^\top\mathbf f_i$。完整流程见 \cref{algo:on-rnea}。

\begin{algo}[RNEA：递推 Newton--Euler 逆动力学]\label{algo:on-rnea}
输入 $\mathbf q,\dot{\mathbf q},\ddot{\mathbf q}$，输出 $\boldsymbol\tau$。初始化 $\mathbf v_0\!\leftarrow\!\mathbf 0$，$\mathbf a_0\!\leftarrow\!-\mathbf g_0$。
\begin{itemize}
  \item \textbf{Pass 1（前向，$i=1..n$）}：$[{}^i\mathbf X_{\lambda(i)},\mathbf S_i]\!\leftarrow\!\mathrm{jcalc}(i,q_i,\dot q_i)$；$\mathbf v_i\!\leftarrow\!{}^i\mathbf X_{\lambda(i)}\mathbf v_{\lambda(i)}+\mathbf S_i\dot q_i$；$\mathbf c_i\!\leftarrow\!\mathbf v_i\times(\mathbf S_i\dot q_i)$；$\mathbf a_i\!\leftarrow\!{}^i\mathbf X_{\lambda(i)}\mathbf a_{\lambda(i)}+\mathbf S_i\ddot q_i+\mathbf c_i$；$\mathbf f_i\!\leftarrow\!\mathcal I_i\mathbf a_i+\mathbf v_i\times^*(\mathcal I_i\mathbf v_i)-\mathbf f_i^{\text{ext}}$。
  \item \textbf{Pass 2（反向，$i=n..1$）}：$\tau_i\!\leftarrow\!\mathbf S_i^\top\mathbf f_i$；若 $\lambda(i)\neq0$ 则 $\mathbf f_{\lambda(i)}\mathrel{+}={}^{\lambda(i)}\mathbf X_i^*\mathbf f_i$。
\end{itemize}
\end{algo}

\begin{theorem}[RNEA 的 $O(n)$ 复杂度]\label{thm:on-rnea-complexity}
RNEA 的计算量严格正比于关节数 $n$，约 $130n$ FLOPs。
\end{theorem}
\begin{proof}[证明]
两个 pass 都是树上的\emph{单次}遍历，无嵌套循环。Pass 1 每关节为常数开销：Plücker 变换 $\approx24$ FLOPs（利用稀疏结构）、运动叉积 $\approx12$、加速度递推 $\approx30$、$\mathcal I_i\mathbf a_i$（对称）$\approx21$、力叉积 $\approx12$，合计 $k_1\approx100$。Pass 2 每关节：$\mathbf S_i^\top\mathbf f_i$（revolute 仅 1 非零）$=1$、力回溯 $\approx24$、向量加 $6$，合计 $k_2\approx30$。总 FLOPs $=n(k_1+k_2)\approx130n$，故 $O(n)$。$\blacksquare$
\end{proof}

\begin{theorem}[RNEA 正确性]\label{thm:on-rnea-correct}
$\mathrm{RNEA}(\mathbf q,\dot{\mathbf q},\ddot{\mathbf q})=\mathbf M(\mathbf q)\ddot{\mathbf q}+\mathbf C(\mathbf q,\dot{\mathbf q})\dot{\mathbf q}+\mathbf g(\mathbf q)$。
\end{theorem}
\begin{proof}[证明骨架]
由 \cref{eq:on-arec}，$\mathbf c_i$ 不含 $\ddot{\mathbf q}$，故 $\mathbf a_i$ 是 $\ddot{\mathbf q}$ 的仿射函数 $\mathbf a_i=\mathbf a_i^{(\ddot q)}+\mathbf a_i^{(0)}$，其中 $\mathbf a_i^{(\ddot q)}=\sum_{j}\boldsymbol\Phi_{ij}\mathbf S_j\ddot q_j$（$\boldsymbol\Phi_{ij}$ 为 link $j\to i$ 的 Plücker 变换链，仅 $j$ 是 $i$ 祖先时非零）。代入 \cref{eq:on-frec} 并投影到 $\mathbf S_i^\top$，提取正比于 $\ddot{\mathbf q}$ 的系数矩阵——它恰是 $\mathbf M(\mathbf q)$ 的 CRBA 定义（Walker--Orin 1982）。\textbf{零阶项}：令 $\dot{\mathbf q}=\ddot{\mathbf q}=\mathbf 0$，则 $\mathbf c_i=\mathbf 0$、$\mathbf v_i=\mathbf 0$、velocity product force 为零，仅剩 $\mathbf a_0=-\mathbf g_0$ 传播产生的力，输出为 $\mathbf g(\mathbf q)$。\textbf{二次项}：剩余「含 $\dot{\mathbf q}$ 不含 $\ddot{\mathbf q}$」的项即 $\mathbf C(\mathbf q,\dot{\mathbf q})\dot{\mathbf q}$（$\mathbf c_i$ 对 $\dot{\mathbf q}$ 一次、velocity product force 对 $\dot{\mathbf q}$ 二次），与 Christoffel 符号给出的 $\mathbf C$ 矩阵一致。$\blacksquare$
\end{proof}

\paragraph{特殊调用。}
RNEA 的通用性在于：设不同输入可提取 Lagrange 方程的各部分（\cref{tab:on-rnea-calls}）。第四种用法尤具启发：跑 $n$ 次 RNEA（每次 $\ddot{\mathbf q}=\mathbf e_j$）即可逐列构出完整 $\mathbf M(\mathbf q)$，复杂度 $O(n^2)$——与 CRBA 同阶但常数大 2--3 倍（这正是 \cref{sec:on-crba} 改进的动机）。

\begin{table}[H]\centering\small
\caption{RNEA 的特殊调用。}
\label{tab:on-rnea-calls}
\begin{tabular}{lll}
\toprule
调用 & 输出 & 用途 \\
\midrule
$\mathrm{RNEA}(\mathbf q,\dot{\mathbf q},\ddot{\mathbf q})$ & $\mathbf M\ddot{\mathbf q}+\mathbf C\dot{\mathbf q}+\mathbf g$ & 完整逆动力学 \\
$\mathrm{RNEA}(\mathbf q,\dot{\mathbf q},\mathbf 0)$ & $\mathbf C(\mathbf q,\dot{\mathbf q})\dot{\mathbf q}+\mathbf g(\mathbf q)$ & 非线性效应（重力+科氏补偿） \\
$\mathrm{RNEA}(\mathbf q,\mathbf 0,\mathbf 0)$ & $\mathbf g(\mathbf q)$ & 纯重力项 \\
$\mathrm{RNEA}(\mathbf q,\mathbf 0,\mathbf e_j)$ & $\mathbf M$ 第 $j$ 列 & 逐列构造质量矩阵 \\
\bottomrule
\end{tabular}
\end{table}

\begin{note}
\textbf{易错：bias acceleration 漏写.}\;
若把 \cref{eq:on-arec} 写成 $\mathbf a_i={}^i\mathbf X_{\lambda(i)}\mathbf a_{\lambda(i)}+\mathbf S_i\ddot q_i$（漏掉 $\mathbf c_i$），则 RNEA 对 $\dot{\mathbf q}=\mathbf 0$ 仍正确、但 $\dot{\mathbf q}\neq\mathbf 0$ 时错误——因为 $\mathbf c_i$ 是科氏加速度，$\ddot q_i=0$ 时也存在。自检：比较 $\mathrm{RNEA}(\mathbf q,\dot{\mathbf q},\ddot{\mathbf q})$ 与 $\mathbf M(\mathbf q)\ddot{\mathbf q}+\mathrm{RNEA}(\mathbf q,\dot{\mathbf q},\mathbf 0)$，差异应小于 $10^{-10}$。
\end{note}

\section{CRBA：质量矩阵 \texorpdfstring{$\mathbf M(\mathbf q)$}{M(q)} 的 \texorpdfstring{$O(n^2)$}{O(n^2)}}
\label{sec:on-crba}

\paragraph{问题与动机。}
CRBA（复合刚体算法）计算对称正定的广义质量矩阵 $\mathbf M(\mathbf q)$。许多算法需要 $\mathbf M$ \emph{本身}而非 $\mathbf M\ddot{\mathbf q}$：全身控制 QP 的约束 $\mathbf M\ddot{\mathbf q}+\mathbf h=\mathbf S^\top\boldsymbol\tau+\mathbf J_c^\top\boldsymbol\lambda$（\cref{ch:force-wbc}）、阻抗控制、状态估计的 $\ddot{\mathbf q}=\mathbf M^{-1}(\boldsymbol\tau-\mathbf h)$、操作空间惯性 $\boldsymbol\Lambda=(\mathbf J\mathbf M^{-1}\mathbf J^\top)^{-1}$（\cref{ch:operational-space}）。RNEA 无法直接给出 $\mathbf M$；逐列调 $n$ 次 RNEA 可行但常数大，CRBA 用一次后向遍历填满 $\mathbf M$。

\paragraph{复合刚体惯性。}
核心概念：把以 link $i$ 为根的子树「冻结」成一个刚体，其总空间惯性即\emph{复合刚体惯性} $\mathcal I_i^c$（即 \cref{eq:sva-Icomposite}）。为什么冻结子树？因为 $M_{ij}$ 的物理含义是「仅关节 $j$ 以单位加速度运动、其余冻结时，关节 $i$ 需要的力矩」；当子树关节全冻结，子树就是一个刚体，惯性可一次算出，无需逐 link 累加。

\begin{insight}
CRBA 的复合刚体惯性本质是\textbf{自底向上的动态规划}：它利用 $\mathbf M(\mathbf q)$ 的深层结构——$M_{ij}$ 只依赖关节 $i,j$ 最低公共祖先到两者路径上的 link 惯性。把子树惯性累加后，路径信息被压缩成 6×6 矩阵，大幅减少重复计算（类比从叶子开始把每个小分枝重量累加到母枝，根处即得全树重量，$O(1)$ 查询）。
\end{insight}

\paragraph{算法与复杂度。}
见 \cref{algo:on-crba}：后向遍历中，对角块 $M_{ii}=\mathbf S_i^\top\mathcal I_i^c\mathbf S_i$，沿祖先链传播 $\mathbf F={}^j\mathbf X_{\lambda(j)}^\top\mathbf F$ 填充交叉块 $M_{ji}=\mathbf S_j^\top\mathbf F$（含义：关节 $i$ 加速时祖先 $j$ 承受的反力）。复杂度 $O(n)+\sum_i\mathrm{depth}(i)$：链式拓扑 $\sum_i i=O(n^2)$；平衡树 $\sum_i\log i=O(n\log n)$。$O(n^2)$ 已是填充 $\mathbf M$ 上三角全部非零元的信息论下界——CRBA 通过共享子树惯性把暴力法的 $O(n^3)$ 降到 $O(n^2)$。

\begin{algo}[CRBA：复合刚体算法（仅填上三角）]\label{algo:on-crba}
输入 $\mathbf q$，输出 $\mathbf M\in\mathbb R^{n\times n}$（上三角）。前向 $i=1..n$：$[{}^i\mathbf X_{\lambda(i)},\mathbf S_i]\!\leftarrow\!\mathrm{jcalc}(i,q_i)$，$\mathcal I_i^c\!\leftarrow\!\mathcal I_i$。后向 $i=n..1$：
\begin{itemize}
  \item 若 $\lambda(i)\neq0$：$\mathcal I_{\lambda(i)}^c\mathrel{+}={}^i\mathbf X_{\lambda(i)}^\top\mathcal I_i^c\,{}^i\mathbf X_{\lambda(i)}$；
  \item $\mathbf F\!\leftarrow\!\mathcal I_i^c\mathbf S_i$；$M_{ii}\!\leftarrow\!\mathbf S_i^\top\mathbf F$；
  \item $j\!\leftarrow\!i$；\textbf{while} $\lambda(j)\neq0$：$\mathbf F\!\leftarrow\!{}^j\mathbf X_{\lambda(j)}^\top\mathbf F$，$j\!\leftarrow\!\lambda(j)$，$M_{ji}\!\leftarrow\!\mathbf S_j^\top\mathbf F$。
\end{itemize}
\end{algo}

\begin{example}[2 连杆臂的 CRBA]\label{ex:on-crba-2link}
设 link 1、2 质量 $m_1=3$、$m_2=2$ kg，两 revolute 关节。后向从 $i=2$：$\mathcal I_2^c=\mathcal I_2$（叶子），$M_{22}=\mathbf S_2^\top\mathcal I_2^c\mathbf S_2$（对 revolute z 即 $\mathcal I_2^c$ 的 $(3,3)$ 元素，link 2 绕关节轴的转动惯量）；将 $\mathcal I_2^c$ 经 Plücker 变换加到 $\mathcal I_1^c$；填 $M_{12}=\mathbf S_1^\top({}^1\mathbf X_2^\top\mathcal I_2^c\mathbf S_2)$（关节 2 加速对关节 1 的反力矩）。再 $i=1$：此时 $\mathcal I_1^c$ 已含 link 2 贡献，$M_{11}=\mathbf S_1^\top\mathcal I_1^c\mathbf S_1$。三处观察：（i）$\mathcal I_1^c$ 含整臂惯性，故根关节有效惯量最大 $M_{11}>M_{22}$；（ii）$M_{12}=M_{21}$，对称性来自动能 $T=\tfrac12\dot{\mathbf q}^\top\mathbf M\dot{\mathbf q}$ 对 $\dot{\mathbf q}$ 的 Hessian 必对称；（iii）结果与 Lagrange 法（对动能取偏导）完全一致。
\end{example}

\begin{note}
\textbf{易错：CRBA 只返回上三角.}\;
多数库（如 Pinocchio）的 CRBA 只填 $\mathbf M$ 上三角以省一半计算，下三角为零或垃圾值。直接 $\mathbf M\ddot{\mathbf q}$ 会错；须先对称化 $\mathbf M\!\leftarrow\!\mathbf M+\mathbf M^\top-\mathrm{diag}(\mathbf M)$，或用「对称视图」求解器告知 $\mathbf M$ 对称。
\end{note}

\section{ABA：正动力学 \texorpdfstring{$O(n)$}{O(n)}}
\label{sec:on-aba}

\paragraph{问题与动机。}
ABA（关节体算法）求解正动力学：已知 $\mathbf q,\dot{\mathbf q},\boldsymbol\tau$，求加速度
\begin{equation}
  \ddot{\mathbf q}=\mathrm{ABA}(\mathbf q,\dot{\mathbf q},\boldsymbol\tau)=\mathbf M(\mathbf q)^{-1}\big(\boldsymbol\tau-\mathbf C(\mathbf q,\dot{\mathbf q})\dot{\mathbf q}-\mathbf g(\mathbf q)\big).
  \label{eq:on-aba-def}
\end{equation}
仿真器的核心循环就是正动力学：给定 $(\mathbf q,\dot{\mathbf q})$ 与 $\boldsymbol\tau$ 求 $\ddot{\mathbf q}$ 再积分。Lagrange 路线（CRBA + Cholesky + 求解）是 $O(n^2)+O(n^3/3)+O(n^2)$；Featherstone（1983）的洞察是：\textbf{不必先构造 $\mathbf M$ 再求逆，可在树上递推中直接消元}，把一切压缩到 $O(n)$\cite{featherstone1983calculation}。

\paragraph{关节体惯性。}
对 link $i$ 及其子树，假设所有子关节\emph{自由}（只受约束力、无外驱力矩），则存在唯一对称正定的 6×6 矩阵 $\mathcal I_i^A$ 与 6D 偏置力 $\mathbf p_i^A$ 使 $\mathbf f_i=\mathcal I_i^A\mathbf a_i+\mathbf p_i^A$。它与复合刚体惯性 $\mathcal I_i^c$ 形成对比（\cref{tab:on-IA-Ic}）。

\begin{table}[H]\centering\small
\caption{复合刚体惯性 $\mathcal I^c$ 与关节体惯性 $\mathcal I^A$ 的对比。}
\label{tab:on-IA-Ic}
\begin{tabular}{lll}
\toprule
特征 & 复合刚体惯性 $\mathcal I_i^c$ & 关节体惯性 $\mathcal I_i^A$ \\
\midrule
假设 & 子树关节\textbf{冻结}（$\ddot{\mathbf q}=\dot{\mathbf q}=\mathbf 0$） & 子树关节\textbf{自由}（$\boldsymbol\tau=\mathbf 0$） \\
含义 & 「冻结子树」的总惯量 & 「铰接子树」的等效惯量 \\
数学性质 & 总 $\succeq\mathcal I_i$ & 总 $\preceq\mathcal I_i^c$（半正定序） \\
用于 & CRBA（算 $\mathbf M$） & ABA（算正动力学） \\
\bottomrule
\end{tabular}
\end{table}

跨领域类比：推一辆购物车，若车轮锁死（冻结）你感受整车惯性 $\mathcal I^c$；若车轮自由转动（铰接），等效惯性更小——车轮转动「吸收」了部分推力。$\mathcal I^A$ 正是「考虑关节自由度后的等效惯性」。

\paragraph{Schur 补递推。}
$\mathcal I^A,\mathbf p^A$ 从叶到根递推，每层消去子关节约束力：
\begin{equation}
  \mathcal I_i^A=\mathcal I_i+\sum_{j\in\text{children}(i)}{}^i\mathbf X_j^*\Big(\mathcal I_j^A-\mathbf U_j\mathbf D_j^{-1}\mathbf U_j^\top\Big){}^j\mathbf X_i,
  \qquad \mathbf U_j=\mathcal I_j^A\mathbf S_j,\;\; \mathbf D_j=\mathbf S_j^\top\mathbf U_j.
  \label{eq:on-aba-schur}
\end{equation}
括号内 $\mathcal I_j^A-\mathbf U_j\mathbf D_j^{-1}\mathbf U_j^\top$ 正是经典的 \textbf{Schur 补}。其几何含义：二次型 $\mathbf x^\top\mathbf A\mathbf x$ 分块为 $\begin{psmallmatrix}\mathbf A_{11}&\mathbf A_{12}\\ \mathbf A_{21}&\mathbf A_{22}\end{psmallmatrix}$，若 $\mathbf x_2$ 自由（可选最优值最小化二次型），对 $\mathbf x_2$ 求最优后剩余 $\mathbf x_1$ 的有效矩阵是 $\mathbf A_{11}-\mathbf A_{12}\mathbf A_{22}^{-1}\mathbf A_{21}$。在 ABA 中，$\mathbf x_2$ 对应自由的关节加速度 $\ddot q_j$、$\mathbf x_1$ 对应传给父节点的加速度——「消去 $\ddot q_j$」即「让关节自由运动，看父节点感受多少等效惯性」。对 revolute 关节 $\mathbf D_j$ 是标量，$\mathbf U_j\mathbf D_j^{-1}\mathbf U_j^\top$ 是秩-1 更新。完整三遍流程见 \cref{algo:on-aba}。

\begin{algo}[ABA：关节体正动力学（三遍）]\label{algo:on-aba}
输入 $\mathbf q,\dot{\mathbf q},\boldsymbol\tau$，输出 $\ddot{\mathbf q}$。
\begin{itemize}
  \item \textbf{Pass 1（前向，$i=1..n$）}：$\mathbf v_i\!\leftarrow\!{}^i\mathbf X_{\lambda(i)}\mathbf v_{\lambda(i)}+\mathbf S_i\dot q_i$；$\mathbf c_i\!\leftarrow\!\mathbf v_i\times(\mathbf S_i\dot q_i)$；$\mathcal I_i^A\!\leftarrow\!\mathcal I_i$；$\mathbf p_i^A\!\leftarrow\!\mathbf v_i\times^*(\mathcal I_i\mathbf v_i)-\mathbf f_i^{\text{ext}}$。
  \item \textbf{Pass 2（反向，$i=n..1$）}：$\mathbf U_i\!\leftarrow\!\mathcal I_i^A\mathbf S_i$；$\mathbf D_i\!\leftarrow\!\mathbf S_i^\top\mathbf U_i$；$u_i\!\leftarrow\!\tau_i-\mathbf S_i^\top\mathbf p_i^A$；若 $\lambda(i)\neq0$：$\mathcal I_i^{A,\mathrm{red}}\!\leftarrow\!\mathcal I_i^A-\mathbf U_i\mathbf D_i^{-1}\mathbf U_i^\top$，$\mathcal I_{\lambda(i)}^A\mathrel{+}={}^i\mathbf X_{\lambda(i)}^\top\mathcal I_i^{A,\mathrm{red}}\,{}^i\mathbf X_{\lambda(i)}$，$\mathbf p_{\lambda(i)}^A\mathrel{+}={}^i\mathbf X_{\lambda(i)}^\top(\mathbf p_i^A+\mathcal I_i^A\mathbf c_i+\mathbf U_i\mathbf D_i^{-1}u_i)$。
  \item \textbf{Pass 3（前向，$i=1..n$）}，$\mathbf a_0\!\leftarrow\!-\mathbf g_0$：$\mathbf a_i'\!\leftarrow\!{}^i\mathbf X_{\lambda(i)}\mathbf a_{\lambda(i)}+\mathbf c_i$；$\ddot q_i\!\leftarrow\!\mathbf D_i^{-1}(u_i-\mathbf U_i^\top\mathbf a_i')$；$\mathbf a_i\!\leftarrow\!\mathbf a_i'+\mathbf S_i\ddot q_i$。
\end{itemize}
\end{algo}

\begin{theorem}[ABA 的复杂度与正确性]\label{thm:on-aba}
ABA 严格 $O(n)$（约 $200n$ FLOPs），且 $\mathrm{ABA}(\mathbf q,\dot{\mathbf q},\boldsymbol\tau)=\mathbf M(\mathbf q)^{-1}(\boldsymbol\tau-\mathbf C\dot{\mathbf q}-\mathbf g)$。
\end{theorem}
\begin{proof}[证明骨架]
\emph{复杂度}：三遍均为树上单次遍历无嵌套循环，每关节常数开销（Pass 1$\approx60$、Pass 2$\approx100$、Pass 3$\approx40$ FLOPs），总 $\approx200n$。\emph{正确性}：Pass 1 与 RNEA 的 Pass 1 同（算 $\mathbf v_i,\mathbf c_i$），$\mathbf p_i^A$ 是「$\ddot{\mathbf q}=\mathbf 0$ 时每 link 施加的 spatial force」。Pass 2 逐关节 Schur 补消元，数学上等价于对线性方程组 $\mathbf M\ddot{\mathbf q}=\boldsymbol\tau-\mathbf C\dot{\mathbf q}-\mathbf g$ 从最后关节起逐个消去 $\ddot q_n,\ddot q_{n-1},\dots$，$\mathcal I_{\lambda(i)}^A$ 记录消元后的等效惯性。Pass 3 是回代——从根向叶逐个解 $\ddot q_i$。由构造，输出满足 $\mathbf M\ddot{\mathbf q}+\mathbf h=\mathbf S^\top\boldsymbol\tau$（无约束时），即 \cref{eq:on-aba-def}。$\blacksquare$
\end{proof}

\begin{theorem}[ABA 等价于块 $\mathbf{LDL^\top}$ 分解]\label{thm:on-aba-ldlt}
ABA 的三遍在代数上等价于对 $\mathbf M(\mathbf q)$ 做\emph{块 $\mathrm{LDL}^\top$ 分解}（块结构由关节树诱导）：Pass 2 = 因子化步骤（从最后行起消元），Pass 3 = 前向/后向替换求解 $\mathbf M\ddot{\mathbf q}=\boldsymbol\tau-\mathbf h$；$\mathbf D_i$ 是 $\mathrm{LDL}^\top$ 对角块的第 $i$ 元素，$\mathbf U_i$ 是 $\mathbf L$ 因子的列块\cite{featherstone2008rigid,jain2011robot}。
\end{theorem}

\begin{insight}
\cref{thm:on-aba-ldlt} 揭示 ABA 不是「天外飞仙」的巧妙算法，而是一个\textbf{自然的稀疏线性代数操作}——它在利用 $\mathbf M(\mathbf q)$ 的树状稀疏结构做高效分解。$\mathbf M$ 的 Cholesky 因子 $\mathbf L$ 的非零结构恰是「每个节点与其所有祖先」（由关节树拓扑决定）。这一视角的三重威力：（i）解释了为何 ABA 是 $O(n)$（树状 $\mathrm{LDL}^\top$ 即 $O(n)$）；（ii）指明推广方向（闭链约束打破树状结构、需处理 fill-in，见 \cref{ch:constrained-dynamics}）；（iii）连接动力学与 SLAM——iSAM2 的 Bayes 树（\cref{ch:isam2}）也是在 elimination tree 上做稀疏因子化，「动力学的稀疏分解 = SLAM 后端的稀疏分解」。
\end{insight}

\begin{example}[2 自由度臂的 ABA 手推]\label{ex:on-aba-2link}
2 自由度平面臂（两 revolute z），$\mathbf q=[0,0]^\top$、$\dot{\mathbf q}=[1,2]^\top$、$\boldsymbol\tau=[5,3]^\top$。\textbf{Pass 1}：$\mathbf v_1=\mathbf S_1=(0,0,1,0,0,0)^\top$，$\mathbf c_1=\mathbf v_1\times(\mathbf S_1\dot q_1)=\mathbf 0$（同方向叉积为零）；$\mathbf v_2={}^2\mathbf X_1\mathbf v_1+\mathbf S_2\dot q_2$，$\mathbf c_2\neq\mathbf 0$（$\mathbf v_2$ 线速度部分非零）；$\mathcal I_i^A\!\leftarrow\!\mathcal I_i$，$\mathbf p_i^A\!\leftarrow\!\mathbf v_i\times^*(\mathcal I_i\mathbf v_i)$。\textbf{Pass 2}（$i=2$ 起）：$\mathbf U_2=\mathcal I_2^A\mathbf S_2$（取第 3 列），$\mathbf D_2=\mathcal I_2^A[3,3]$（标量），$u_2=3-\mathbf p_2^A[3]$，Schur 补 $\mathcal I_2^{A,\mathrm{red}}=\mathcal I_2^A-\mathbf U_2\mathbf D_2^{-1}\mathbf U_2^\top$（秩-1 更新）经变换加到 $\mathcal I_1^A$；再 $i=1$：$\mathbf D_1$ 已含子树贡献。\textbf{Pass 3}：$\mathbf a_0=-\mathbf g_0$，$\ddot q_1=\mathbf D_1^{-1}(u_1-\mathbf U_1^\top\mathbf a_1')$，$\ddot q_2=\mathbf D_2^{-1}(u_2-\mathbf U_2^\top\mathbf a_2')$。全程未构造 $\mathbf M$ 也未求逆，直接从 $\boldsymbol\tau$ 得 $\ddot{\mathbf q}$。
\end{example}

\begin{note}
\textbf{易错：$\mathcal I^A$ 正定性丢失.}\;
Schur 补 $\mathcal I^A-\mathbf U\mathbf D^{-1}\mathbf U^\top$ 在浮点下可能失去对称正定性，尤其当 $\mathbf D_i$ 很小（关节有效惯量近零，如细长轻质末端）时数值误差被放大、$\mathbf D_i^{-1}\to\infty$、输出爆炸。工程修复：给 $\mathbf D_i$ 加 armature（电机反射惯量，$I_{\text{armature}}=I_{\text{rotor}}r^2$，$r$ 为减速比）使 $\mathbf D_i\geq I_{\text{armature}}>0$；或显式对称化 $\mathcal I^A\!\leftarrow\!(\mathcal I^A+\mathcal I^{A\top})/2$；或用双精度。注意 ABA 在精确算术下与 $\mathbf M^{-1}(\boldsymbol\tau-\mathbf h)$ 完全等价，差异仅浮点级（$\sim10^{-12}$）。
\end{note}

\section{Minv 与三算法统一视角}
\label{sec:on-unify}

\paragraph{Minv 直接算法。}
某些场景需完整 $\mathbf M^{-1}$ 矩阵（如操作空间惯性、接触求解的 Delassus 矩阵 $\mathbf G=\mathbf J\mathbf M^{-1}\mathbf J^\top$）。关键洞察：$\mathbf M^{-1}$ 第 $j$ 列即 $\mathbf M^{-1}\mathbf e_j=\mathrm{ABA}(\mathbf q,\mathbf 0,\mathbf e_j)$（对「仅关节 $j$ 有单位力矩、不注入重力」做正动力学）。简单调 $n$ 次 ABA 是 $O(n^2)$，但 Minv 直接算法更聪明：$n$ 次调用共享同一 Pass 1+2（$\mathcal I^A,\mathbf p^A,\mathbf U,\mathbf D$ 不依赖 $\boldsymbol\tau$），只需 $n$ 次不同的 Pass 3，总 $O(n)+n\cdot O(n)=O(n^2)$、常数比 CRBA+Cholesky+求逆更小。而若只需 $\mathbf M^{-1}\mathbf v$（单个向量），直接 $\mathrm{ABA}(\mathbf q,\mathbf 0,\mathbf v)$ 即 $O(n)$——这是 MPC 中 100 个 horizon 节点上 $100\times O(n)$ vs $100\times O(n^2)$ 的差距。

\paragraph{消息传递统一。}
RNEA、CRBA、ABA 看似功能不同，本质都是\emph{同一棵关节树上的消息传递}（\cref{tab:on-unify}），并与概率图模型的推断算法惊人地平行：RNEA $\leftrightarrow$ Belief Propagation（树上两次消息传递）、CRBA $\leftrightarrow$ Variable Elimination（自底向上消元）、ABA $\leftrightarrow$ Bayes 网上的高斯消元（$\mathrm{LDL}^\top$ 因子化+回代）。

\begin{table}[H]\centering\small
\caption{三大算法的消息传递统一视角。}
\label{tab:on-unify}
\begin{tabular}{llll}
\toprule
算法 & 功能 & 消息方向 & 传递的「消息」 \\
\midrule
RNEA & 逆动力学 $\boldsymbol\tau$ & 前向+后向 & 速度/加速度 $\to$ 力/力矩 \\
CRBA & 质量矩阵 $\mathbf M$ & 仅后向 & 复合惯性 $\mathcal I^c$ \\
ABA & 正动力学 $\ddot{\mathbf q}$ & 前向+后向+前向 & 速度 $\to$ 关节体惯性 $\mathcal I^A$ $\to$ 加速度 \\
\bottomrule
\end{tabular}
\end{table}

\begin{insight}
三算法的深层统一源自 $\mathbf M(\mathbf q)$ 的稀疏 Cholesky 因子恰好是关节树的 elimination tree。RNEA 是「矩阵-向量乘」的 $O(n)$ 实现（隐式利用 $\mathbf M$ 的因子化），CRBA 是「显式构造 $\mathbf M$」，ABA 是「直接求解 $\mathbf M\mathbf x=\mathbf b$ 而不显式构造 $\mathbf M$」。Jain 的空间算子代数（SOA）给出最优雅的表达 $\mathbf M=\mathbf H\boldsymbol\Phi\mathbf M_0\boldsymbol\Phi^\top\mathbf H^\top$\cite{jain2011robot}（$\boldsymbol\Phi$ 为树上传播算子=Plücker 变换链、$\mathbf H$ 为关节选择矩阵=$\mathbf S$、$\mathbf M_0$ 为块对角 link 惯性）：RNEA 计算 $\mathbf M\ddot{\mathbf q}+\mathbf h$，CRBA 显式构造该乘积，ABA 利用其 $\mathrm{LDL}^\top$ 因子化求解。
\end{insight}

\paragraph{Lagrange vs Newton--Euler 的计算粒度。}
二者物理等价、给出相同运动方程，但计算效率有本质差异（\cref{tab:on-lag-ne}）。根源：Lagrange 法把「局部关节耦合」编码成「全局矩阵元素」，即使 $\mathbf M$ 多数元素为零，构造矩阵本身也需 $O(n^2)$；Newton--Euler 递推从不构造该矩阵，直接在树上传「力消息」、每 link 只与邻居通信。这与 PageRank 同理：可先构造 $O(n^2)$ 邻接矩阵再做矩阵乘，或直接在图上做 $O(\text{边数})$ 消息传递。

\begin{table}[H]\centering\small
\caption{Lagrange 法与 Newton--Euler 递推的计算粒度对比。}
\label{tab:on-lag-ne}
\begin{tabular}{lll}
\toprule
维度 & Lagrange & Newton--Euler（递推） \\
\midrule
思考方式 & 能量（标量）$\to$ 变分 $\to$ 运动方程 & 力/力矩（向量）$\to$ 牛顿定律 $\to$ 递推 \\
计算粒度 & 整个系统的能量表达式 & 每 link 的局部力平衡 \\
利用拓扑 & 困难（$\mathbf M$ 稠密） & 自然（树递推直接用稀疏性） \\
最优复杂度 & $O(n^2)$ (CRBA) / $O(n^3)$ (直接) & $O(n)$ (RNEA/ABA) \\
\bottomrule
\end{tabular}
\end{table}

\paragraph{选型经验。}
逆动力学唯一选 RNEA；需 $\mathbf M$ 用 CRBA；正动力学小自由度（$n\lesssim12$）CRBA+密集 Cholesky 反而更快（BLAS/SIMD 对密集小矩阵极度优化、ABA 在 $n=7$ 时甚至慢于 CRBA+Cholesky），$n\gtrsim12$ 用 ABA、$n=36$ 时 ABA 快约 4 倍（交叉点经验值 $n\approx9$\cite{featherstone2008rigid}）；只需 $\mathbf M^{-1}\mathbf v$ 用 $\mathrm{ABA}(\mathbf q,\mathbf 0,\mathbf v)$。

\section{稀疏结构与并行（理论部分）}
\label{sec:on-sparse}

\paragraph{Branch-induced sparsity。}
当自由度从串联臂（7）扩到人形（36+）或多机器人（100+），$\mathbf M(\mathbf q)$ 维度激增，但其非零结构由关节树诱导、高度稀疏。

\begin{theorem}[分支诱导稀疏性]\label{thm:on-sparsity}
$M_{ij}\neq0$ 当且仅当关节 $i$ 与 $j$ 在关节树中具有祖先--后代关系\cite{featherstone2005empirical}。等价地，$\mathbf M$ 的非零结构恰是关节树的「可达性矩阵」。
\end{theorem}

不同拓扑的稀疏性差异巨大：串联链上三角全非零、无 fill-in，因子化 $O(n^2)$；平衡二叉树（人形双腿双臂）每行约 $O(\log n)$ 非零，因子化约 $O(n\log^2 n)$；星形（多腿共享躯干）仅躯干行全满，因子化 $O(n)+O(k^2)$（$k$ 为躯干自由度）。

\paragraph{稀疏 $\mathrm{LDL}^\top$ 与 elimination tree。}
稀疏 Cholesky 的实现核心是：外层循环从 $n-1$ 到 $0$（子树后于父节点处理）、内层只\emph{沿祖先链上爬}填充 $\mathbf U$ 的非零元——这是稀疏性的代码体现。链式拓扑祖先链长 $O(n)$ 故总 $O(n^2)$，平衡树祖先链长 $O(\log n)$ 故 $O(n\log n)$。这一结构与 SLAM 的 iSAM2 共享同一数学对象 elimination tree（\cref{tab:on-elim}）：两者都在树结构诱导的稀疏线性系统上做高效因子化。理解动力学的稀疏 Cholesky，即理解 SLAM 后端优化的核心数据结构——这是机器人学中「动力学」与「估计」共享深层数学结构的绝佳例证，双向交叉引用至 \cref{ch:isam2}。

\begin{table}[H]\centering\small
\caption{动力学稀疏分解 $\leftrightarrow$ SLAM 后端的共同抽象。}
\label{tab:on-elim}
\begin{tabular}{lll}
\toprule
动力学 & SLAM & 共同抽象 \\
\midrule
关节树 & 因子图 & 图结构 \\
$\mathbf M(\mathbf q)$ 的稀疏模式 & 信息矩阵的稀疏模式 & 稀疏矩阵 \\
CRBA 的祖先链遍历 & 变量消元的 clique tree & elimination tree \\
ABA / 稀疏 $\mathrm{LDL}^\top$ & iSAM2 的 Bayes 树更新 & 增量因子化 \\
\bottomrule
\end{tabular}
\end{table}

\paragraph{D\&C ABA 的 $O(\log n)$ 思想。}
标准 RNEA/ABA 是\emph{串行}算法（前向第 $i$ 步依赖第 $i-1$ 步）。Featherstone 的分治 ABA（IJRR 1999）\cite{featherstone1999divide} 把 $n$ 关节树从中间切成两个 $n/2$ 子树，各自独立计算关节体惯性 $\mathcal I^A$，在连接点做一次 6×6 Schur 补合并（合并操作 $O(1)$，不影响并行深度）。给定 $O(n)$ 个处理器，递归深度 $O(\log n)$、每层 $O(1)$ 个 6×6 操作，故\textbf{总并行时间 $O(\log n)$}。主要用于超大多体系统（如分子动力学 1000+ 链节）；对常规 36 自由度机器人，$O(n)$ 已足够快（约 3$\,\mu\mathrm s$），分治的通信开销与负载不均（人形树不平衡）不划算。

\section{本章脉络与展望}
\label{sec:on-outlook}

\paragraph{全景。}
本章沿一条主线把刚体动力学的计算地基铺到底：\textbf{6D 代数工具}（twist/wrench、对偶功率、$\mathbf X/\mathbf X^*$、$\times/\times^*$、空间惯性 $\mathcal I$）$\to$ \textbf{李代数统一}（$\mathrm{ad}/\mathrm{ad}^*/\mathrm{Ad}$，$\mathfrak{se}(3)$）$\to$ \textbf{$O(n)$ 递推算法}（RNEA 逆动力学、CRBA 质量矩阵、ABA 正动力学）$\to$ \textbf{稀疏 $\mathrm{LDL}^\top$ 与并行}（elimination tree、D\&C ABA）。\cref{fig:sva-roadmap} 把这条主线绘成全景图。其中两个最深的理论高地——「Coriolis 隐藏在力叉积里」（\cref{eq:sva-NE}）与「ABA 等价于块 $\mathrm{LDL}^\top$」（\cref{thm:on-aba-ldlt}）——分别接通几何力学的 Christoffel 视角与 SLAM 后端的 Bayes 树。

\begin{figure}[ht]
\centering
\begin{tikzpicture}[
  node distance=5mm and 8mm, >=Stealth,
  blk/.style={draw, rounded corners, align=center, font=\small, inner sep=4pt, fill=main!6, draw=main!60!black},
  grp/.style={draw, rounded corners, align=center, font=\small, inner sep=4pt, fill=second!8, draw=second!60!black},
  ln/.style={->, thick, gray!60!black}]
  \node[blk] (tools) {6D 代数工具\\ twist/wrench, $\mathbf X/\mathbf X^*$\\ $\times/\times^*$, $\mathcal I$};
  \node[blk, right=of tools] (lie) {李代数统一\\ $\mathrm{ad},\mathrm{ad}^*,\mathrm{Ad}$\\ $\cong\mathfrak{se}(3)$};
  \node[grp, below=8mm of tools] (rnea) {RNEA $O(n)$\\ 逆动力学 $\boldsymbol\tau$};
  \node[grp, right=of rnea] (crba) {CRBA $O(n^2)$\\ 质量矩阵 $\mathbf M$};
  \node[grp, right=of crba] (aba) {ABA $O(n)$\\ 正动力学 $\ddot{\mathbf q}$};
  \node[blk, below=8mm of crba, fill=third!8, draw=third!60!black] (sparse) {稀疏 $\mathrm{LDL}^\top$ / 并行\\ elimination tree, D\&C};
  \draw[ln] (tools)--(lie);
  \draw[ln] (tools.south) -- ++(0,-3mm) -| (rnea.north);
  \draw[ln] (lie.south) -- ++(0,-3mm) -| (aba.north);
  \draw[ln] (rnea)--(crba); \draw[ln] (crba)--(aba);
  \draw[ln] (aba.south) -- ++(0,-3mm) -| (sparse.north);
\end{tikzpicture}
\caption{本章主线：从 6D 代数工具到 $O(n)$ 递推算法再到稀疏分解。}
\label{fig:sva-roadmap}
\end{figure}

\paragraph{启下：兄弟章与跨部钩子。}
本章是刚体动力学子部的地基，向多个方向延伸：
\begin{itemize}
  \item \textbf{几何力学（\cref{ch:geometric-mechanics}）}：本章给「Coriolis 在 $\hat{\mathbf v}\times^*\mathcal I\hat{\mathbf v}$ 中」的递推视角，几何力学给 Lagrange/Hamilton 变分来源、Christoffel 符号、SE(3) 几何力学与辛积分的解析视角；$\mathbf M,\mathbf C,\mathbf g$ 的定义与性质主放该章，本章用结论。
  \item \textbf{约束动力学（\cref{ch:constrained-dynamics}）}：本章 ABA 的 Schur 补与稀疏 $\mathrm{LDL}^\top$ 是闭链/接触约束（约束力 $\boldsymbol\lambda\to$ KKT/Delassus $\mathbf G=\mathbf J\mathbf M^{-1}\mathbf J^\top$）的前置；解析微分（$\partial\mathrm{RNEA}/\partial\mathbf q$、Carpentier--Mansard RSS 2018\cite{carpentier2018analytical}）主体留该章，求导扰动采本书右扰动约定（\cref{ch:lie}）。
  \item \textbf{操作空间与力控（\cref{ch:operational-space,ch:force-wbc}）}：本章 $\mathbf M^{-1}$ 是操作空间惯性 $\boldsymbol\Lambda=(\mathbf J\mathbf M^{-1}\mathbf J^\top)^{-1}$（Khatib 1987\cite{khatib1987unified}）与 WBC/TSID 约束 $\mathbf M\ddot{\mathbf q}+\mathbf h=\mathbf S^\top\boldsymbol\tau+\mathbf J_c^\top\boldsymbol\lambda$ 的计算地基。
  \item \textbf{接触力学（\cref{ch:contact-mechanics}）}：接触力作为额外空间力 $\mathbf f_i^{\text{ext}}$ 注入 RNEA 力回溯（\cref{eq:on-frec}），N 章讲摩擦锥/GIWC/LCP 主体。
  \item \textbf{控制部 DDP/iLQR 与 MPC（\cref{ch:ddp-ilqr,ch:mpc-advanced}）}：每节点需 ABA（正动力学）+ 解析导数；本章算法是其动力学内核。机械臂部、腿足部把本章通用 $O(n)$ 算法特化到固定基串联臂与浮基系统（浮基用虚拟 6 自由度关节，$n_q\neq n_v$，基座 RNEA 返回值前 6 维是 base wrench 而非驱动力）。
\end{itemize}

\section*{本章小结}
\addcontentsline{toc}{section}{本章小结}

本章把空间向量代数与 $O(n)$ 动力学递推算法编织成一条从「6D 几何工具」到「微秒级实时代码地基」的连续主线。核心结论汇于 \cref{tab:sva-summary}。

\begin{table}[H]\centering\small
\caption{本章核心公式与算法速查。}
\label{tab:sva-summary}
\begin{tabular}{lll}
\toprule
对象 & 核心式 & 出处 \\
\midrule
twist / wrench & $\hat{\mathbf v}=(\boldsymbol\omega;\mathbf v_O)$ / $\hat{\mathbf f}=(\boldsymbol\tau_O;\mathbf f)$ & \cref{eq:sva-twist,eq:sva-wrench} \\
功率对偶 & $P=\hat{\mathbf f}^\top\hat{\mathbf v}$（参考点无关） & \cref{eq:sva-power} \\
运动变换 / 力变换 & ${}^B\mathbf X_A=\begin{psmallmatrix}\mathbf R&\mathbf 0\\ [\mathbf p]_\times\mathbf R&\mathbf R\end{psmallmatrix}$, $\mathbf X^*=\mathbf X^{-\top}$ & \cref{eq:sva-X,eq:sva-Xstar} \\
运动叉积 / 力叉积 & $[\hat{\mathbf v}]_\times$（下三角）, $[\hat{\mathbf v}]_{\times^*}=-[\hat{\mathbf v}]_\times^\top$（上三角） & \cref{eq:sva-crm,eq:sva-crf} \\
空间惯性 & $\mathcal I_O$（10 参数，对称）, $\mathcal I_B=\mathbf X^{-\top}\mathcal I_A\mathbf X^{-1}$ & \cref{eq:sva-IO,eq:sva-Iparallel} \\
单体 Newton--Euler & $\hat{\mathbf f}=\mathcal I\hat{\mathbf a}+\hat{\mathbf v}\times^*(\mathcal I\hat{\mathbf v})$ & \cref{eq:sva-NE} \\
李代数统一 & $\mathbf X=\mathrm{Ad}_{\mathbf T}$, $[\hat{\mathbf v}]_\times=\mathrm{ad}$, $\mathbf X^*=\mathrm{Ad}^*$, $[\hat{\mathbf v}]_{\times^*}=\mathrm{ad}^*$ & \cref{thm:sva-lie} \\
RNEA / CRBA / ABA & 逆动力学 $O(n)$ / 质量矩阵 $O(n^2)$ / 正动力学 $O(n)$ & \cref{algo:on-rnea,algo:on-crba,algo:on-aba} \\
ABA $\equiv\mathrm{LDL}^\top$ & 块 $\mathrm{LDL}^\top$ 分解（树诱导稀疏） & \cref{thm:on-aba-ldlt} \\
\bottomrule
\end{tabular}
\end{table}

\section*{常见误解辨析}
\addcontentsline{toc}{section}{常见误解辨析}

\begin{itemize}
  \item \textbf{「空间向量就是把 3D 向量拼成 6D」}：错。空间向量有严格代数结构——生活在互相对偶的两空间中，遵循不同变换规则（运动用 $\mathbf X$、力用 $\mathbf X^{-\top}$），有李代数结构（运动叉积=李括号）、内蕴度量（空间惯性）。威力不在「简化记号」而在「6D 空间的代数闭合性」。
  \item \textbf{「空间加速度是 link 质心的线加速度加角加速度」}：错。空间加速度 $=\mathrm d(\text{空间速度})/\mathrm dt$，其线性部分是 $\dot{\mathbf v}_O$ 而非质心加速度 $\ddot{\mathbf r}$，二者差 $\boldsymbol\omega\times\mathbf v$（\cref{eq:sva-accel-cls}）。基座注入 $\mathbf a_0=-\mathbf g_0$ 用的正是空间加速度约定。
  \item \textbf{「Coriolis 力在空间框架中消失了」}：错。它没消失，而是隐藏在力叉积 $\hat{\mathbf v}\times^*(\mathcal I\hat{\mathbf v})$ 中（\cref{eq:sva-NE}）；RNEA 力回溯的该项投影到关节空间求和即 $\mathbf C(\mathbf q,\dot{\mathbf q})\dot{\mathbf q}$。
  \item \textbf{「RNEA 只能做逆动力学」}：错。其输出含 $\mathbf M\ddot{\mathbf q}+\mathbf C\dot{\mathbf q}+\mathbf g$ 整体信息，经特殊调用（\cref{tab:on-rnea-calls}）可提取 $\mathbf g$、$\mathbf C\dot{\mathbf q}+\mathbf g$、$\mathbf M$ 各列。
  \item \textbf{「ABA 是近似算法，不如 $\mathbf M^{-1}(\boldsymbol\tau-\mathbf h)$ 精确」}：错。ABA 在精确算术下与之代数等价，浮点差异仅 $\sim10^{-12}$（双精度）。
  \item \textbf{「Drake 也用 6D 速度，可直接套 Pinocchio 的 6×6 变换公式」}：错。Drake 明确声明其空间向量「不是 Plücker 向量」，用显式 \texttt{Shift}/\allowbreak\texttt{Rotate} 而非 6×6 矩阵乘，直接套用会得错误结果（\cref{tab:sva-notations}）。
\end{itemize}

\section*{延伸阅读}
\addcontentsline{toc}{section}{延伸阅读}

空间向量代数与三大算法的权威主教材是 Featherstone 的 \emph{Rigid Body Dynamics Algorithms}（2008）第 2、5、6、7 章\cite{featherstone2008rigid}；面向初学者的直觉性导引见其 IEEE RAM 2010 两部分综述\cite{featherstone2010beginner}。三大算法的原始论文分别是 RNEA（Luh--Walker--Paul 1980\cite{luh1980online}）、CRBA（Walker--Orin 1982）、ABA（Featherstone 1983\cite{featherstone1983calculation}）；递推 Lagrange 路线见 Hollerbach 1980\cite{hollerbach1980recursive} 与等价性证明 Silver 1982\cite{silver1982equivalence}。李群视角的 twist/wrench 理论见 Murray--Li--Sastry 1994\cite{murray1994mathematical} 与 Lynch--Park \emph{Modern Robotics}\cite{lynch2017modern}；空间算子代数的统一框架见 Jain 2011\cite{jain2011robot}。稀疏 $\mathrm{LDL}^\top$ 见 Featherstone IJRR 2005\cite{featherstone2005empirical}，并行 D\&C ABA 见 Featherstone IJRR 1999\cite{featherstone1999divide}，解析微分见 Carpentier--Mansard RSS 2018\cite{carpentier2018analytical}，现代高性能实现见 Pinocchio\cite{carpentier2019pinocchio}。物理一致性的三角不等式见 Wensing 等\cite{wensing2018linear}；螺旋分解的 Mozzi 优先权见 Ceccarelli\cite{ceccarelli2000screw}，螺旋代数用于空间机构见 Yang--Freudenstein 1964\cite{yang1964application}。
